\providecommand{\SO}{\operatorname{SO}}
\providecommand{\OO}{\operatorname{O}}
\providecommand{\SE}{\operatorname{SE}}
\providecommand{\Eucl}{\operatorname{E}}
\providecommand{\GL}{\operatorname{GL}}
\providecommand{\Sym}{\mathfrak{S}}
\providecommand{\Haar}{\mathrm{d}\mu}
\providecommand{\softmax}{\operatorname{softmax}}
\providecommand{\Stab}{\operatorname{Stab}}
\providecommand{\Orb}{\operatorname{Orb}}
\providecommand{\Attn}{\operatorname{Attn}}
\providecommand{\MLP}{\operatorname{MLP}}
\providecommand{\diag}{\operatorname{diag}}

\chapter{Grupos: Redes Equivariantes a Grupos}
\label{ch:groups}

Las redes convolucionales del \Cref{ch:grids} deben su \'exito a una sola
simetr\'ia: la traslaci\'on. Un filtro que reconoce un borde en una esquina de
una imagen reconoce el mismo borde en cualquier otra parte, porque la
convoluci\'on conmuta con los desplazamientos. Pero la traslaci\'on es solo una
simetr\'ia entre muchas. Una imagen de tejido al microscopio no tiene un
``arriba'' privilegiado; una fotograf\'ia de sat\'elite puede aparecer rotada un
\'angulo cualquiera; una mol\'ecula puede reflejarse en su imagen especular. Una
red equivariante solo a traslaciones debe redescubrir de forma independiente cada
patr\'on rotado o reflejado, malgastando capacidad en copias de la misma
caracter\'istica.

Este cap\'itulo desarrolla la respuesta sistem\'atica: reemplazar el grupo de
traslaciones por un grupo de simetr\'ia arbitrario $\group$ y generalizar la
convoluci\'on en consecuencia. El resultado ---la \emph{convoluci\'on de grupo}---
produce redes equivariantes por construcci\'on a rotaciones, reflexiones y otras
transformaciones geom\'etricas. Construimos la teor\'ia a partir de grupos finitos
como las simetr\'ias del cuadrado, la extendemos a grupos compactos continuos
($\SO(2)$, $\SO(3)$) mediante filtros orientables y an\'alisis arm\'onico, la
llevamos a los grupos euclidianos no compactos ($\Eucl(3)$, $\SE(3)$) que rigen
las nubes de puntos y las mol\'eculas tridimensionales, y cerramos con el grupo de
permutaciones, la simetr\'ia de los conjuntos no ordenados. A lo largo del
cap\'itulo permanece vigente el principio rector establecido en
\Cref{ch:priors}: las capas equivariantes se componen en redes equivariantes, y
una \'unica lectura invariante al final produce una predicci\'on insensible a la
pose elegida.

\section{M\'as all\'a de las traslaciones}
\label{sec:groups:beyond}

Recu\'erdese del \Cref{ch:priors} que una aplicaci\'on $\Phi$ entre espacios de
se\~nales es \emph{$\group$-equivariante} si $\Phi(L_g f) = L'_g\,\Phi(f)$ para
todo $g \in \group$, donde $L_g$ y $L'_g$ denotan la acci\'on de $g$ sobre los
espacios de entrada y de salida, e \emph{$\group$-invariante} si $\Phi(L_g f) =
\Phi(f)$. La equivarianza generaliza la invarianza traslacional de las redes
convolucionales cl\'asicas: si los datos poseen simetr\'ias bajo $\group$,
entonces las aplicaciones $\group$-equivariantes preservan esas simetr\'ias a lo
largo de todo el procesamiento, dando lugar a representaciones m\'as robustas y
m\'as eficientes en par\'ametros. El lenguaje matem\'atico de grupos, acciones de
grupo y representaciones se desarrolla en \Cref{ch:foundations}; aqu\'i lo
ponemos a trabajar para construir arquitecturas.

La red convolucional cl\'asica es el caso particular $\group = (\Z^2, +)$ que
act\'ua por traslaciones sobre im\'agenes discretas. En este caso las
representaciones irreducibles son unidimensionales (los caracteres $x \mapsto
e^{2\pi i \langle k, x\rangle}$), la convoluci\'on de grupo se reduce a la
convoluci\'on ordinaria y los filtros son invariantes a traslaciones. Las
extensiones geom\'etricas de este cap\'itulo incorporan grupos m\'as ricos:
$\SO(2)$ para rotaciones planas, $\SE(2)$ para movimientos r\'igidos planos,
$\SO(3)$ y $\SE(3)$ para rotaciones y movimientos r\'igidos tridimensionales, y
el grupo sim\'etrico $\Sym_n$ para permutaciones de colecciones no ordenadas. La
receta de dise\~no es siempre la misma:
\begin{enumerate}
  \item \textbf{Identificar el grupo}: determinar qu\'e simetr\'ias est\'an
        presentes en los datos.
  \item \textbf{Analizar las representaciones}: caracterizar las
        representaciones irreducibles relevantes bajo las cuales transforman las
        caracter\'isticas.
  \item \textbf{Dise\~nar la convoluci\'on}: implementar operaciones
        equivariantes de manera eficiente usando la estructura de grupo.
  \item \textbf{Optimizar}: explotar la simetr\'ia para compartir par\'ametros y
        reducir el espacio de hip\'otesis sin p\'erdida de poder expresivo.
\end{enumerate}

\section{Convoluciones de grupo}
\label{sec:groups:convolutions}

\subsection{Convoluci\'on sobre un grupo}
\label{sec:groups:conv-on-group}

La generalizaci\'on de la convoluci\'on a un grupo arbitrario es el n\'ucleo
te\'orico de las extensiones geom\'etricas de las redes convolucionales. Los
fundamentos de la teor\'ia de representaciones de grupos en que se apoya esta
construcci\'on se remontan al trabajo cl\'asico de \citet{Peter1927} y
\citet{Serre1977}.

\begin{definition}[Convoluci\'on de grupo]
\label{def:groups:group-convolution}
Para funciones $f, \psi: \group \to \R$ definidas sobre un grupo $\group$, la
\emph{convoluci\'on de grupo} es
\begin{equation}
  (f *_{\group} \psi)(u)
  = \int_{\group} f(uv^{-1})\,\psi(v)\,\Haar(v),
  \label{eq:groups:group-convolution}
\end{equation}
donde $\mu$ es la medida de Haar sobre $\group$.
\end{definition}

Esta definici\'on generaliza la convoluci\'on ordinaria reemplazando la resta
$u - v$ por la operaci\'on de grupo $uv^{-1}$. Para grupos discretos la integral
se convierte en una suma.

El siguiente teorema establece la profunda conexi\'on entre equivarianza y
convoluci\'on: la convoluci\'on no es meramente \emph{una} forma de construir una
capa equivariante, sino esencialmente la \emph{\'unica} forma.

\begin{theorem}[Kondor--Trivedi]
\label{thm:groups:kondor-trivedi}
Una red neuronal feedforward $N$ es equivariante a la acci\'on de un grupo
compacto $\group$ sobre sus entradas si y solo si cada capa de $N$ implementa una
forma generalizada de convoluci\'on derivada de
\Cref{eq:groups:group-convolution}
\citep{kondor2018generalizationconvolutionalneural}.
\end{theorem}

\begin{proof}[Esbozo]
La demostraci\'on tiene dos direcciones. \emph{Directa:} si cada capa implementa
una convoluci\'on de grupo, entonces por construcci\'on preserva la estructura de
grupo, lo que implica equivarianza (verificado directamente m\'as abajo en
\Cref{thm:groups:lift-equivariance,thm:groups:group-conv-equivariance}).
\emph{Rec\'iproca:} si la red es equivariante, entonces la teor\'ia de
representaciones caracteriza todas las aplicaciones lineales equivariantes
admisibles como convoluciones generalizadas. El argumento completo usa el
an\'alisis arm\'onico sobre $\group$ y la descomposici\'on de Peter--Weyl para
describir toda aplicaci\'on lineal equivariante entre espacios de
caracter\'isticas.
\end{proof}

La elecci\'on de la representaci\'on bajo la cual transforman las
caracter\'isticas determina qu\'e tipos de patrones puede detectar una capa de
manera equivariante. En una red $\group$-equivariante, cada capa procesa
caracter\'isticas que transforman bajo una representaci\'on espec\'ifica $\rho:
\group \to \GL(V)$ del grupo de simetr\'ia.

\begin{example}[Representaci\'on regular en redes convolucionales de grupo]
\label{ex:groups:regular-representation}
Para un grupo finito $\group = \{g_1, \dots, g_n\}$, la \emph{representaci\'on
regular izquierda} act\'ua sobre funciones $f: \group \to \R^d$ mediante
\begin{equation}
  (\rho(g) f)(h) = f(g^{-1} h).
\end{equation}
Esta representaci\'on es fundamental en las redes convolucionales de grupo, donde
los canales corresponden a las evaluaciones de la funci\'on en los distintos
elementos del grupo.
\end{example}

Por el teorema de Peter--Weyl (\Cref{ch:foundations}), el espacio $L^2(\group)$
de un grupo compacto se descompone como $L^2(\group) = \bigoplus_{\lambda \in
\widehat \group} n_\lambda \mathcal{H}_\lambda$, una suma directa sobre las
representaciones irreducibles $\mathcal{H}_\lambda$ con multiplicidades
$n_\lambda$. Esto tiene una consecuencia directa para el dise\~no de
arquitecturas: cada $\mathcal{H}_\lambda$ corresponde a un tipo de canal, y
$n_\lambda$ prescribe cu\'antos canales de ese tipo incluir para capturar por
completo las simetr\'ias del grupo.

En la pr\'actica, las se\~nales no est\'an definidas sobre todo el grupo sino
sobre un \emph{espacio homog\'eneo} $X = \group / H$, donde $H$ es un subgrupo.

\begin{definition}[Espacio homog\'eneo]
\label{def:groups:homogeneous-space}
Un espacio $X$ es \emph{homog\'eneo} bajo la acci\'on de $\group$ si para
cualesquiera $x, y \in X$ existe $g \in \group$ con $g \cdot x = y$.
\end{definition}

La esfera $S^2$ es homog\'enea bajo $\SO(3)$, y las im\'agenes ordinarias son
homog\'eneas bajo traslaciones. Generalizar la convoluci\'on a un espacio
homog\'eneo requiere maquinaria adicional de \emph{elevaci\'on} (\emph{lifting})
y \emph{proyecci\'on}, que desarrollamos a continuaci\'on.

\subsection{Capas de elevaci\'on y de convoluci\'on de grupo}
\label{sec:groups:lifting}

Una red convolucional de grupo procede en tres etapas: una \emph{capa de
elevaci\'on} aplica una se\~nal plana a una funci\'on sobre el grupo, una pila de
\emph{capas de convoluci\'on de grupo} procesa esa funci\'on de manera
equivariante, y una \emph{lectura invariante} colapsa la dimensi\'on del grupo en
una predicci\'on.

\paragraph{Elevaci\'on.}
La primera operaci\'on transforma una se\~nal sobre el espacio eucl\'ideo $\R^2$
en una funci\'on sobre el grupo $\group$.

\begin{definition}[Convoluci\'on de elevaci\'on]
\label{def:groups:lifting-convolution}
Dada una funci\'on $f: \R^2 \to \R$ y un filtro $\psi: \R^2 \to \R$, la
\emph{convoluci\'on de elevaci\'on} produce una funci\'on $F: \group \to \R$,
\begin{equation}
  (f \star \psi)(g) = \int_{\R^2} f(x)\,\psi(g^{-1} \cdot x)\,dx,
  \label{eq:groups:lifting-conv}
\end{equation}
donde $g \in \group$ y $g^{-1} \cdot x$ denota la acci\'on de $g^{-1}$ sobre el
punto $x \in \R^2$.
\end{definition}

Para un grupo finito $\group$ la integral es una suma sobre una ventana local, de
modo que para cada posici\'on espacial y cada transformaci\'on $g \in \group$
calculamos la respuesta del filtro $\psi$ transformado por $g^{-1}$.

\begin{theorem}[Equivarianza de la elevaci\'on]
\label{thm:groups:lift-equivariance}
La convoluci\'on de elevaci\'on es equivariante en el sentido de que
\begin{equation}
  (L_h f \star \psi)(g) = (f \star \psi)(h^{-1} g),
\end{equation}
donde $L_h f(x) = f(h^{-1} \cdot x)$ es la acci\'on de $\group$ sobre la entrada.
\end{theorem}

\begin{proof}
Por definici\'on,
\begin{align}
  (L_h f \star \psi)(g)
  &= \int_{\R^2} (L_h f)(x)\,\psi(g^{-1} \cdot x)\,dx
   = \int_{\R^2} f(h^{-1} \cdot x)\,\psi(g^{-1} \cdot x)\,dx.
\end{align}
Con el cambio de variables $y = h^{-1} \cdot x$ (cuyo jacobiano es constante
porque la acci\'on es por isometr\'ias),
\begin{align}
  &= \int_{\R^2} f(y)\,\psi(g^{-1} \cdot h \cdot y)\,dy
   = \int_{\R^2} f(y)\,\psi\big((h^{-1} g)^{-1} \cdot y\big)\,dy
   = (f \star \psi)(h^{-1} g). \qedhere
\end{align}
\end{proof}

\paragraph{Convoluci\'on de grupo.}
Las capas siguientes operan entre funciones sobre el grupo.

\begin{definition}[Capa de convoluci\'on de grupo]
\label{def:groups:group-conv-layer}
Para funciones $f, \psi: \group \to \R$, la \emph{convoluci\'on de grupo} es
\begin{equation}
  (f \star_{\group} \psi)(u) = \sum_{g \in \group} f(g)\,\psi(g^{-1} u),
  \label{eq:groups:group-conv-discrete}
\end{equation}
donde la suma recorre todos los elementos de $\group$. Esta es la contraparte
discreta de \Cref{eq:groups:group-convolution}: el cambio de variables
$v \mapsto g = uv^{-1}$ convierte $\int_{\group} f(uv^{-1})\psi(v)\,\Haar(v)$ en
$\sum_g f(g)\psi(g^{-1}u)$.
\end{definition}

\begin{theorem}[Equivarianza de la convoluci\'on de grupo]
\label{thm:groups:group-conv-equivariance}
La convoluci\'on de grupo es equivariante:
\begin{equation}
  (L_h f) \star_{\group} \psi = L_h (f \star_{\group} \psi),
\end{equation}
donde $L_h$ es la acci\'on de traslaci\'on izquierda $(L_h f)(g) = f(h^{-1} g)$.
\end{theorem}

\begin{proof}
Por definici\'on,
\begin{align}
  \big((L_h f) \star_{\group} \psi\big)(u)
  &= \sum_{g \in \group} (L_h f)(g)\,\psi(g^{-1} u)
   = \sum_{g \in \group} f(h^{-1} g)\,\psi(g^{-1} u).
\end{align}
Con la sustituci\'on $g' = h^{-1} g$ (equivalentemente $g = h g'$),
\begin{align}
  &= \sum_{g' \in \group} f(g')\,\psi\big((h g')^{-1} u\big)
   = \sum_{g' \in \group} f(g')\,\psi\big(g'^{-1} h^{-1} u\big) \\
  &= (f \star_{\group} \psi)(h^{-1} u)
   = \big(L_h (f \star_{\group} \psi)\big)(u). \qedhere
\end{align}
\end{proof}

La convoluci\'on de grupo hereda el \'algebra del \'algebra de grupo
$\R[\group]$, en la que la multiplicaci\'on es la convoluci\'on; esta estructura
algebraica es lo que garantiza que componer capas preserve la equivarianza.

\begin{algorithm}[h]
\caption{Capa de convoluci\'on de grupo para un grupo finito}
\label{alg:groups:gconv-finite}
\begin{algorithmic}[1]
\Require Funci\'on de entrada $f$, filtros $\{\psi_i\}$, grupo
  $\group = \{g_1, \dots, g_n\}$
\Ensure Mapa de caracter\'isticas $\mathrm{output}$
\State Inicializar $\mathrm{output}$ con forma $(|\group|, H, W, C_{\text{out}})$
\For{cada filtro $i \in \{1, \dots, C_{\text{out}}\}$}
  \For{cada elemento de grupo $u \in \group$}
    \For{cada posici\'on $(h, w)$}
      \State $s \leftarrow 0$
      \For{cada elemento de grupo $g \in \group$}
        \For{cada canal de entrada $c$}
          \State $s \leftarrow s + f[g][h][w][c] \cdot \psi_i[u^{-1} g][c]$
        \EndFor
      \EndFor
      \State $\mathrm{output}[u][h][w][i] \leftarrow s$
    \EndFor
  \EndFor
\EndFor
\State \Return $\mathrm{output}$
\end{algorithmic}
\end{algorithm}

\paragraph{Un ejemplo trabajado: el grupo diedral $D_4$.}
El grupo diedral $D_4$ de simetr\'ias del cuadrado es el grupo finito
paradigm\'atico para datos de imagen.

\begin{definition}[El grupo $D_4$]
\label{def:groups:d4}
El grupo diedral $D_4$ tiene ocho elementos,
\begin{equation}
  D_4 = \{e, r, r^2, r^3, s, sr, sr^2, sr^3\},
\end{equation}
donde $e$ es la identidad, $r$ es una rotaci\'on de $90^\circ$ ($r^4 = e$), $s$
es una reflexi\'on ($s^2 = e$), y se cumple la relaci\'on $s r s = r^{-1}$.
\end{definition}

\begin{proposition}[Acci\'on de $D_4$ sobre el plano]
\label{prop:groups:d4-action}
La acci\'on natural de $D_4$ sobre $\R^2$ est\'a generada por
\begin{align}
  r \cdot (x,y) &= (-y, x) && \text{(rotaci\'on de $90^\circ$),} \\
  s \cdot (x,y) &= (x, -y) && \text{(reflexi\'on horizontal).}
\end{align}
Todos los dem\'as elementos surgen por composici\'on.
\end{proposition}

\begin{proof}
La rotaci\'on $r$ act\'ua mediante $R_{90} = \left(\begin{smallmatrix} 0 & -1 \\
1 & 0 \end{smallmatrix}\right)$, dando $R_{90}(x,y)^\top = (-y,x)^\top$. La
reflexi\'on $s$ respecto al eje $x$ act\'ua mediante $S =
\left(\begin{smallmatrix} 1 & 0 \\ 0 & -1 \end{smallmatrix}\right)$, dando
$(x,-y)^\top$. Los seis elementos restantes se siguen por composici\'on: $r^2
\cdot (x,y) = (-x,-y)$, $r^3 \cdot (x,y) = (y,-x)$, $sr \cdot (x,y) = (-y,-x)$,
$sr^2 \cdot (x,y) = (-x,y)$, $sr^3 \cdot (x,y) = (y,x)$. Estas ocho aplicaciones
son exactamente las isometr\'ias lineales de $\R^2$ que preservan el cuadrado
unidad.
\end{proof}

Los ocho elementos satisfacen los axiomas de grupo (clausura, asociatividad,
identidad e inversos), lo que puede verificarse construyendo la tabla de
multiplicaci\'on de $D_4$. \Cref{fig:groups:d4-actions} muestra las ocho
transformaciones actuando sobre un parche asim\'etrico.

\begin{figure}[h]
\centering
\begin{tikzpicture}[
    filled/.style={minimum size=0.55cm, inner sep=0pt, outer sep=0pt,
                   fill=gdlblue!40, draw=gdlblue!60},
    empty/.style={minimum size=0.55cm, inner sep=0pt, outer sep=0pt,
                  fill=gdlgray!15, draw=gdlgray!40},
    grouplabel/.style={font=\large, text=gdlpurple!70!black},
    rowlabel/.style={font=\small\itshape, text=gdlgray!80!black, rotate=90,
                     anchor=south}
]
\def\cs{0.55}
\def\hs{3.0}
\def\vs{3.2}
% TOP ROW: rotations
\begin{scope}[shift={(0, 0)}]
    \node[filled] at (0,0){};\node[filled] at (\cs,0){};\node[filled] at (2*\cs,0){};
    \node[filled] at (0,-\cs){};\node[empty] at (\cs,-\cs){};\node[empty] at (2*\cs,-\cs){};
    \node[filled] at (0,-2*\cs){};\node[filled] at (\cs,-2*\cs){};\node[empty] at (2*\cs,-2*\cs){};
    \node[grouplabel] at (\cs, -2*\cs - 0.7) {$e$};
\end{scope}
\begin{scope}[shift={(\hs, 0)}]
    \node[filled] at (0,0){};\node[empty] at (\cs,0){};\node[empty] at (2*\cs,0){};
    \node[filled] at (0,-\cs){};\node[empty] at (\cs,-\cs){};\node[filled] at (2*\cs,-\cs){};
    \node[filled] at (0,-2*\cs){};\node[filled] at (\cs,-2*\cs){};\node[filled] at (2*\cs,-2*\cs){};
    \node[grouplabel] at (\cs, -2*\cs - 0.7) {$r$};
\end{scope}
\begin{scope}[shift={(2*\hs, 0)}]
    \node[empty] at (0,0){};\node[filled] at (\cs,0){};\node[filled] at (2*\cs,0){};
    \node[empty] at (0,-\cs){};\node[empty] at (\cs,-\cs){};\node[filled] at (2*\cs,-\cs){};
    \node[filled] at (0,-2*\cs){};\node[filled] at (\cs,-2*\cs){};\node[filled] at (2*\cs,-2*\cs){};
    \node[grouplabel] at (\cs, -2*\cs - 0.7) {$r^2$};
\end{scope}
\begin{scope}[shift={(3*\hs, 0)}]
    \node[filled] at (0,0){};\node[filled] at (\cs,0){};\node[filled] at (2*\cs,0){};
    \node[filled] at (0,-\cs){};\node[empty] at (\cs,-\cs){};\node[filled] at (2*\cs,-\cs){};
    \node[empty] at (0,-2*\cs){};\node[empty] at (\cs,-2*\cs){};\node[filled] at (2*\cs,-2*\cs){};
    \node[grouplabel] at (\cs, -2*\cs - 0.7) {$r^3$};
\end{scope}
\node[rowlabel] at (-1.1, -\cs) {Rotaciones};
% BOTTOM ROW: reflections
\begin{scope}[shift={(0, -\vs)}]
    \node[filled] at (0,0){};\node[filled] at (\cs,0){};\node[filled] at (2*\cs,0){};
    \node[empty] at (0,-\cs){};\node[empty] at (\cs,-\cs){};\node[filled] at (2*\cs,-\cs){};
    \node[empty] at (0,-2*\cs){};\node[filled] at (\cs,-2*\cs){};\node[filled] at (2*\cs,-2*\cs){};
    \node[grouplabel] at (\cs, -2*\cs - 0.7) {$s$};
\end{scope}
\begin{scope}[shift={(\hs, -\vs)}]
    \node[empty] at (0,0){};\node[empty] at (\cs,0){};\node[filled] at (2*\cs,0){};
    \node[filled] at (0,-\cs){};\node[empty] at (\cs,-\cs){};\node[filled] at (2*\cs,-\cs){};
    \node[filled] at (0,-2*\cs){};\node[filled] at (\cs,-2*\cs){};\node[filled] at (2*\cs,-2*\cs){};
    \node[grouplabel] at (\cs, -2*\cs - 0.7) {$sr$};
\end{scope}
\begin{scope}[shift={(2*\hs, -\vs)}]
    \node[filled] at (0,0){};\node[filled] at (\cs,0){};\node[empty] at (2*\cs,0){};
    \node[filled] at (0,-\cs){};\node[empty] at (\cs,-\cs){};\node[empty] at (2*\cs,-\cs){};
    \node[filled] at (0,-2*\cs){};\node[filled] at (\cs,-2*\cs){};\node[filled] at (2*\cs,-2*\cs){};
    \node[grouplabel] at (\cs, -2*\cs - 0.7) {$sr^2$};
\end{scope}
\begin{scope}[shift={(3*\hs, -\vs)}]
    \node[filled] at (0,0){};\node[filled] at (\cs,0){};\node[filled] at (2*\cs,0){};
    \node[filled] at (0,-\cs){};\node[empty] at (\cs,-\cs){};\node[filled] at (2*\cs,-\cs){};
    \node[filled] at (0,-2*\cs){};\node[empty] at (\cs,-2*\cs){};\node[empty] at (2*\cs,-2*\cs){};
    \node[grouplabel] at (\cs, -2*\cs - 0.7) {$sr^3$};
\end{scope}
\node[rowlabel] at (-1.1, -\vs - \cs) {Reflexiones};
\end{tikzpicture}
\caption[Elementos de $D_4$ actuando sobre un parche]{Los ocho elementos del
grupo diedral $D_4$ actuando sobre un parche asim\'etrico. Fila superior: las
cuatro rotaciones ($e$, $r$, $r^2$, $r^3$). Fila inferior: las cuatro
reflexiones ($s$, $sr$, $sr^2$, $sr^3$). Una red $D_4$-equivariante produce la
misma salida sin importar qu\'e transformaci\'on se aplique a la entrada.}
\label{fig:groups:d4-actions}
\end{figure}

Una red $D_4$-equivariante completa compone una capa de elevaci\'on, varias capas
de convoluci\'on de grupo y una capa invariante final.

\begin{definition}[Red $D_4$-equivariante]
\label{def:groups:d4-network}
Una red $\Phi: \mathcal{F}(\R^2) \to \R^C$ es \emph{$D_4$-equivariante} si se
descompone como
\begin{equation}
  \Phi = \Phi_{\text{inv}} \circ \Phi_{\group} \circ \Phi_{\group}
         \circ \Phi_{\text{lift}},
\end{equation}
donde $\Phi_{\text{lift}}: \mathcal{F}(\R^2) \to \mathcal{F}(D_4 \times \R^2)$ es
la capa de elevaci\'on, $\Phi_{\group}: \mathcal{F}(D_4 \times \R^2) \to
\mathcal{F}(D_4 \times \R^2)$ son las capas de convoluci\'on de grupo, y
$\Phi_{\text{inv}}: \mathcal{F}(D_4 \times \R^2) \to \R^C$ es una capa invariante
final.
\end{definition}

\begin{theorem}[Composici\'on hacia la invarianza]
\label{thm:groups:composition-invariance}
Si cada componente $\Phi_i$ de \Cref{def:groups:d4-network} es
$D_4$-equivariante y la capa final $\Phi_{\text{inv}}$ es $D_4$-invariante,
entonces la composici\'on $\Phi$ es $D_4$-invariante:
\begin{equation}
  \Phi(L_g f) = \Phi(f) \qquad \forall g \in D_4,\ f \in \mathcal{F}(\R^2).
\end{equation}
\end{theorem}

\begin{proof}
La composici\'on de aplicaciones equivariantes es equivariante (una propiedad
establecida en \Cref{ch:priors}), de modo que $\Phi_{\group} \circ \Phi_{\group}
\circ \Phi_{\text{lift}}$ es equivariante. La capa invariante final elimina
entonces la dependencia del grupo, produciendo una salida invariante.
\end{proof}

En cada capa equivariante las caracter\'isticas se organizan como funciones $f:
D_4 \times \R^2 \to \R^d$; la dimensi\'on del grupo ($|D_4| = 8$) introduce un
factor multiplicativo en el espacio de caracter\'isticas que preserva toda la
informaci\'on sobre orientaciones y reflexiones.

\paragraph{Verificaci\'on de la equivarianza.}
La correcci\'on te\'orica debe complementarse con una medida num\'erica de cu\'an
bien se preserva la equivarianza.

\begin{definition}[Error de equivarianza]
\label{def:groups:equivariance-error}
Para una aplicaci\'on $\Phi$ y un grupo $\group$, el \emph{error de equivarianza}
en una entrada $x$ es
\begin{equation}
  \mathcal{E}_{\group}(\Phi, x)
  = \frac{1}{|\group|} \sum_{g \in \group}
    \norm{\Phi(g \cdot x) - g \cdot \Phi(x)}^2,
\end{equation}
donde $\norm{\cdot}$ es una norma adecuada en el espacio de salida.
\end{definition}

\begin{theorem}[Criterio de equivarianza]
\label{thm:groups:equivariance-criterion}
Una aplicaci\'on $\Phi$ es $\group$-equivariante si y solo si
$\mathcal{E}_{\group}(\Phi,x) = 0$ para toda entrada $x$.
\end{theorem}

\begin{proof}
($\Rightarrow$) Si $\Phi$ es equivariante entonces $\Phi(g \cdot x) = g \cdot
\Phi(x)$ para todo $g$, de modo que cada sumando se anula y
$\mathcal{E}_{\group}(\Phi,x) = 0$. ($\Leftarrow$) Si el promedio de los
t\'erminos no negativos es cero, entonces cada t\'ermino se anula, de modo que
$\Phi(g \cdot x) = g \cdot \Phi(x)$ para todo $g$ por la definici\'on positiva de
la norma.
\end{proof}

\begin{proposition}[Propagaci\'on del error]
\label{prop:groups:error-propagation}
Para una composici\'on $\Phi = \Phi_n \circ \cdots \circ \Phi_1$ en la que cada
$\Phi_i$ tiene constante de Lipschitz $L_i$ y $x_i = (\Phi_{i-1} \circ \cdots
\circ \Phi_1)(x)$ son las activaciones intermedias (con $x_0 = x$), la
desviaci\'on de equivarianza en el peor caso se propaga como
\begin{equation}
  \sup_{g \in \group} \norm{\Phi(g \cdot x) - g \cdot \Phi(x)}
  \leq \sum_{i=1}^{n} \Big(\prod_{j=i+1}^{n} L_j\Big)
    \sup_{g \in \group} \norm{\Phi_i(g \cdot x_i) - g \cdot \Phi_i(x_i)},
\end{equation}
donde el producto vac\'io (para $i = n$) vale $1$.
\end{proposition}

\begin{proof}
Para $n = 1$ la cota es trivial. Para el paso inductivo sea $\Psi =
\Phi_{n-1} \circ \cdots \circ \Phi_1$. Fijando $g \in \group$,
\begin{align}
  \norm{\Phi(g \cdot x) - g \cdot \Phi(x)}
  &= \norm{\Phi_n(\Psi(g \cdot x)) - g \cdot \Phi_n(\Psi(x))} \\
  &\leq \norm{\Phi_n(\Psi(g \cdot x)) - \Phi_n(g \cdot \Psi(x))}
      + \norm{\Phi_n(g \cdot \Psi(x)) - g \cdot \Phi_n(\Psi(x))}.
\end{align}
El primer t\'ermino est\'a acotado por $L_n \norm{\Psi(g \cdot x) - g \cdot
\Psi(x)}$ mediante la continuidad de Lipschitz; el segundo es la desviaci\'on de
equivarianza de $\Phi_n$ en $x_n = \Psi(x)$. Tomando el supremo sobre $g$ y
aplicando la hip\'otesis inductiva a $\Psi$ se obtiene la afirmaci\'on.
\end{proof}

\begin{remark}[Estabilidad num\'erica de la equivarianza]
\label{rem:groups:numerical-stability}
Para una arquitectura $D_4$-equivariante implementada en aritm\'etica de
precisi\'on finita, el error de equivarianza es de orden
\begin{equation}
  \sup_{g \in \group} \norm{\Phi(g \cdot x) - g \cdot \Phi(x)}
  = O\!\big(\epsilon_{\text{mach}} \cdot \operatorname{depth}(\Phi)
    \cdot \norm{x}\big),
\end{equation}
donde $\epsilon_{\text{mach}}$ es la precisi\'on de m\'aquina y
$\operatorname{depth}(\Phi)$ la profundidad de la red. Cada capa introduce un
error de redondeo proporcional a $\epsilon_{\text{mach}}$, y estos se acumulan
linealmente, reflejando el an\'alisis est\'andar de propagaci\'on de errores de
la aritm\'etica de punto flotante. Los errores de orden $10^{-6}$ observados en
la pr\'actica son coherentes con la precisi\'on de punto flotante de 32 bits: la
equivarianza te\'orica se preserva hasta el l\'imite del hardware, y
$\epsilon_{\text{mach}} \to 0$ recupera la equivarianza exacta.
\end{remark}

\subsection{Por qu\'e compensa la equivarianza}
\label{sec:groups:why-pays}

Un resultado central justifica el uso de las convoluciones de grupo: las
arquitecturas equivariantes pueden aproximar cualquier funci\'on equivariante
continua, de modo que la simetr\'ia restringe el espacio de hip\'otesis sin
sacrificar poder expresivo.

\begin{theorem}[Universalidad equivariante]
\label{thm:groups:universality}
Sea $\group$ un grupo compacto que act\'ua sobre un dominio compacto
$\mathcal{D}$, y sea $\mathcal{F}_{\group}$ el espacio de aplicaciones
$\group$-equivariantes continuas $L^2(\mathcal{D}) \to L^2(\mathcal{D})$. Para
toda $f \in \mathcal{F}_{\group}$ y todo $\epsilon > 0$ existe una red
convolucional $\group$-equivariante $\Phi_{\text{GCNN}}$ con $\norm{f -
\Phi_{\text{GCNN}}}_{L^2} < \epsilon$. Es decir, las redes convolucionales de
grupo son aproximadores universales en el espacio de funciones equivariantes.
\end{theorem}

\begin{proof}[Esbozo para $\group = \Z^2$]
Para una $f: \Z^2 \to \R$ de soporte compacto, la transformada de Fourier
discreta da $f(x) = \sum_{k} \hat f(k)\, e^{2\pi i \langle k, x\rangle / N}$.
Cada componente de frecuencia se aproxima mediante una convoluci\'on con un
filtro apropiado, y una red convolucional con $M$ filtros en la primera capa
puede escribirse $\Phi_{\text{CNN}}(f) = \sigma\big(\sum_{j=1}^M w_j (f * h_j) +
b\big)$ para filtros aprendibles $h_j$, pesos $w_j$, sesgo $b$ y activaci\'on
$\sigma$. La equivarianza se cumple autom\'aticamente porque $L_g(f * h) = (L_g
f) * h$, de modo que $\Phi_{\text{CNN}}(L_g f) = L_g \Phi_{\text{CNN}}(f)$. Para
grupos no abelianos el argumento se traslada reemplazando la transformada de
Fourier por el an\'alisis arm\'onico sobre $\group$: la descomposici\'on en
representaciones irreducibles, la convoluci\'on de grupo y la teor\'ia de
representaciones garantizan en conjunto que la red reproduce toda aplicaci\'on
lineal equivariante.
\end{proof}

\begin{corollary}[Eficiencia en par\'ametros]
\label{cor:groups:parameter-efficiency}
Una red $\group$-equivariante requiere $|\group|$ veces menos par\'ametros que
una red sin restricciones para aproximar una funci\'on $\group$-equivariante con
la misma precisi\'on, donde $|\group|$ es el orden del grupo.
\end{corollary}

\begin{proof}
Sea $f: X \to Y$ $\group$-equivariante. Una red sin restricciones debe aprender
el valor de $f$ en cada punto $x$ y en todas sus transformadas $\{g \cdot x : g
\in \group\}$ de manera independiente, requiriendo $N$ par\'ametros para alcanzar
una precisi\'on $\epsilon$. Una red $\group$-equivariante satisface $\Phi(g \cdot
x) = g \cdot \Phi(x)$ por construcci\'on, de modo que basta aprender $\Phi$ en un
dominio fundamental de la acci\'on, de tama\~no $|X|/|\group|$; los valores en
las restantes $|\group| - 1$ copias se siguen por equivarianza. Esto reduce el
n\'umero de par\'ametros a $N/|\group|$ para la misma precisi\'on.
\end{proof}

Esto precisa por qu\'e las arquitecturas equivariantes generalizan mejor:
explotan autom\'aticamente las simetr\'ias de los datos, encogiendo el espacio de
hip\'otesis sin p\'erdida de poder expresivo.

\paragraph{Grupos finitos en visi\'on.}
Los grupos finitos m\'as relevantes para el procesamiento de im\'agenes preservan
la rejilla discreta de p\'ixeles a la vez que codifican simetr\'ias
significativas: el grupo c\'iclico $C_4 = \{e, r, r^2, r^3\}$ de rotaciones de
$90^\circ$, el grupo diedral $D_4$ que a\~nade reflexiones, y los grupos m\'as
peque\~nos $C_2$ (simetr\'ia bilateral) y $D_2$. Estos grupos preservan la
estructura de rejilla, aparecen de forma natural en texturas, cristales e
im\'agenes m\'edicas, y admiten implementaciones exactas por ser finitos. Para el
grupo c\'iclico $C_4$, una red equivariante requiere exactamente una cuarta parte
de los par\'ametros de una red est\'andar que representa el mismo espacio de
funciones $C_4$-equivariantes: las cuatro copias rotadas $\{\psi, r\cdot\psi,
r^2\cdot\psi, r^3\cdot\psi\}$ se generan a partir de un \'unico filtro base
$\psi$ por la acci\'on del grupo, una instancia directa de
\Cref{cor:groups:parameter-efficiency}.

\begin{remark}[Compromiso dimensionalidad--expresividad]
\label{rem:groups:tradeoff}
Para una red convolucional de grupo que procesa funciones sobre $\group \times
\Omega$, la dimensi\'on del espacio de caracter\'isticas escala como $|\group|
\cdot \dim(\Omega) \cdot C$. Sin embargo, el reparto de pesos inducido por la
equivarianza reduce el n\'umero de par\'ametros libres. Si los datos poseen
simetr\'ias adicionales ---de modo que el estabilizador $\Stab(x) = \{g \in
\group : g \cdot x = x\}$ es no trivial para las entradas t\'ipicas--- la
dimensi\'on efectiva de par\'ametros se reduce a\'un m\'as. Por el teorema
\'orbita--estabilizador, $|\group| = |\Orb(x)| \cdot |\Stab(x)|$, lo que sugiere
de forma heur\'istica que la dimensi\'on efectiva es de orden $|\Orb(x)| \cdot
\dim(\Omega) \cdot C$ en lugar de $|\group| \cdot \dim(\Omega) \cdot C$.
Computacionalmente, el coste crece linealmente con el orden del grupo,
$\text{FLOPs}_{\text{GCNN}} = |\group| \cdot \text{FLOPs}_{\text{CNN}}$, lo que
puede ser prohibitivo para grupos grandes.
\end{remark}

En benchmarks rotados el beneficio es dr\'astico. Los experimentos de
\citet{cohen2016group} muestran que sobre una versi\'on rotada de CIFAR-10 una
red convolucional est\'andar cae a aproximadamente un $47\%$ de precisi\'on,
mientras que las redes $C_4$- y $D_4$-equivariantes conservan en torno al $72\%$
y el $73\%$, con una ligera mejora tambi\'en en la tarea sin rotar. El precio es
la sensibilidad a las desviaciones de la simetr\'ia exacta: en datos reales la
simetr\'ia suele ser solo aproximada, de modo que el grupo debe elegirse con
cuidado.

\section{Filtros orientables y representaciones irreducibles}
\label{sec:groups:steerable}

Los grupos finitos capturan exactamente simetr\'ias discretas, pero una red
$C_4$-equivariante no puede representar un objeto rotado $45^\circ$: cualquier
orientaci\'on intermedia se aproxima por la m\'as cercana discreta, con un error
de \emph{aliasing} que puede alcanzar $45^\circ$ para $C_4$. El remedio es pasar
al grupo continuo de rotaciones $\SO(2)$ y representar los filtros en una base
que rota anal\'iticamente.

\subsection{Filtros orientables}
\label{sec:groups:steerable-filters}

\begin{definition}[El grupo $\SO(2)$]
\label{def:groups:so2}
El grupo ortogonal especial $\SO(2) = \{R_\theta : \theta \in [0, 2\pi)\}$ consta
de todas las rotaciones planas en torno al origen, con
\begin{equation}
  R_\theta = \begin{pmatrix} \cos\theta & -\sin\theta \\
                             \sin\theta & \cos\theta \end{pmatrix}.
\end{equation}
\end{definition}

\begin{proposition}[Propiedades de $\SO(2)$]
\label{prop:groups:so2-properties}
El grupo $\SO(2)$ es \emph{compacto} (topol\'ogicamente un c\'irculo $S^1$),
\emph{abeliano} ($R_\alpha R_\beta = R_{\alpha+\beta} = R_\beta R_\alpha$) y
\emph{uniparam\'etrico} (parametrizado por un \'unico \'angulo).
\end{proposition}

\begin{proof}
\emph{Compacto:} la aplicaci\'on $\theta \mapsto R_\theta$ es un homeomorfismo de
$[0,2\pi)/(0 \sim 2\pi)$ sobre $\SO(2)$, identific\'andolo con $S^1$, que es
cerrado y acotado en $\R^{2 \times 2}$, y por tanto compacto por
Heine--Borel. \emph{Abeliano:} $R_\alpha R_\beta$ es la rotaci\'on por $\alpha +
\beta$, y la adici\'on m\'odulo $2\pi$ es conmutativa. \emph{Uniparam\'etrico:}
las condiciones $R^\top R = I$ y $\det R = 1$ imponen tres restricciones
independientes sobre las cuatro entradas de una matriz $2 \times 2$, dejando $4 -
3 = 1$ grado de libertad.
\end{proof}

La estructura circular hace del an\'alisis arm\'onico la herramienta natural. Un
\emph{filtro orientable} es aquel cuya rotaci\'on puede sintetizarse a partir de
una base finita.

\begin{definition}[Filtro orientable]
\label{def:groups:steerable-filter}
Un filtro $h: \R^2 \to \R$ es \emph{orientable} de orden $N$ si existen filtros
base $\{h_0, \dots, h_{N-1}\}$ y funciones de interpolaci\'on $\{k_n(\theta)\}$
tales que el filtro rotado $h_\theta$ satisface
\begin{equation}
  h_\theta(x,y) = \sum_{n=0}^{N-1} k_n(\theta)\, h_n(x,y).
\end{equation}
\end{definition}

Los arm\'onicos circulares proporcionan una base orientable natural, porque
transforman mediante una simple fase bajo rotaci\'on.

\begin{theorem}[Representaci\'on en arm\'onicos circulares]
\label{thm:groups:circular-harmonics}
Toda funci\'on $f: S^1 \to \C$ admite una serie de Fourier $f(\theta) = \sum_{k
\in \Z} c_k e^{ik\theta}$. Bajo una rotaci\'on, $g_\phi(\theta) = f(\theta +
\phi)$ tiene coeficientes $c_k e^{ik\phi}$.
\end{theorem}

\begin{proof}
El desarrollo se sigue de la completitud del sistema ortonormal
$\{e^{ik\theta}\}_{k \in \Z}$ en $L^2(S^1)$, con $c_k = \frac{1}{2\pi}
\int_0^{2\pi} f(\theta) e^{-ik\theta}\,d\theta$. Para la transformaci\'on,
\[
  \hat g_k = \frac{1}{2\pi}\int_0^{2\pi} f(\theta + \phi) e^{-ik\theta}\,d\theta
  = \frac{1}{2\pi}\int_0^{2\pi} f(\alpha) e^{-ik(\alpha - \phi)}\,d\alpha
  = e^{ik\phi} c_k,
\]
usando la sustituci\'on $\alpha = \theta + \phi$ y la periodicidad.
\end{proof}

\begin{definition}[Filtro arm\'onico circular]
\label{def:groups:circular-harmonic-filter}
Un filtro arm\'onico circular de frecuencia $k$ en coordenadas polares $(r,
\theta)$ es $\psi_k(r, \theta) = R_k(r)\, e^{ik\theta}$, con perfil radial
$R_k(r)$ y parte angular $e^{ik\theta}$.
\end{definition}

\Cref{fig:groups:steerable} ilustra la orientabilidad de un filtro de borde
orientado: cualquier rotaci\'on del filtro base es una combinaci\'on lineal fija
de dos filtros base.

\begin{figure}[h]
\centering
\begin{tikzpicture}[
    >=Stealth,
    every node/.style={font=\small},
    section label/.style={font=\normalsize\bfseries, text=gdlgray!70!black}
]
\newcommand{\steerablefilter}[4]{%
    \begin{scope}[shift={(#1,#2)}, rotate=#3]
        \fill[white, draw=gdlblue!60, thick] (0,0) -- (90:#4) arc (90:270:#4) -- cycle;
        \fill[gdlblue!30, draw=gdlblue!60, thick] (0,0) -- (270:#4) arc (270:450:#4) -- cycle;
        \draw[gdlblue!60, thick] (0,0) circle (#4);
        \draw[gdlgray!50, thin] (0,-#4) -- (0,#4);
    \end{scope}
}
\newcommand{\steerablefiltercolor}[5]{%
    \begin{scope}[shift={(#1,#2)}, rotate=#3]
        \fill[white, draw=#5, thick] (0,0) -- (90:#4) arc (90:270:#4) -- cycle;
        \fill[gdlblue!30, draw=#5, thick] (0,0) -- (270:#4) arc (270:450:#4) -- cycle;
        \draw[#5, thick] (0,0) circle (#4);
        \draw[gdlgray!50, thin] (0,-#4) -- (0,#4);
    \end{scope}
}
\node[section label, anchor=west] at (-1.2, 2.8) {Rotaciones del filtro base};
\steerablefilter{0}{1.2}{0}{0.7}
\node[below, text=gdlblue!70!black, font=\small] at (0, 0.35) {Filtro base $\psi_0$};
\draw[->, thick, gdlgray!60] (1.0, 1.2) -- (1.8, 1.2);
\steerablefilter{2.8}{1.2}{0}{0.6}
\node[below, text=gdlpurple!70!black, font=\scriptsize] at (2.8, 0.45) {$0^\circ$};
\steerablefilter{4.6}{1.2}{45}{0.6}
\node[below, text=gdlpurple!70!black, font=\scriptsize] at (4.6, 0.45) {$45^\circ$};
\steerablefilter{6.4}{1.2}{90}{0.6}
\node[below, text=gdlpurple!70!black, font=\scriptsize] at (6.4, 0.45) {$90^\circ$};
\steerablefilter{8.2}{1.2}{135}{0.6}
\node[below, text=gdlpurple!70!black, font=\scriptsize] at (8.2, 0.45) {$135^\circ$};
\draw[gdlpurple!60, thick, decorate, decoration={brace, amplitude=5pt, mirror}]
    (2.1, 0.1) -- (8.9, 0.1);
\node[below=6pt, text=gdlpurple!70!black, font=\small] at (5.5, 0.05) {$\rho(r_\theta)\psi_0$};
\draw[->, gdlorange!70, thick, dashed] (3.5, 2.0) to[bend left=25] (4.5, 2.0);
\draw[->, gdlorange!70, thick, dashed] (5.3, 2.0) to[bend left=25] (6.3, 2.0);
\draw[->, gdlorange!70, thick, dashed] (7.1, 2.0) to[bend left=25] (8.1, 2.0);
\node[above, font=\tiny, text=gdlorange!70!black] at (4.0, 2.2) {$+45^\circ$};
\node[above, font=\tiny, text=gdlorange!70!black] at (5.8, 2.2) {$+45^\circ$};
\node[above, font=\tiny, text=gdlorange!70!black] at (7.6, 2.2) {$+45^\circ$};
\draw[gdlgray!30, thin] (-1.2, -0.7) -- (10.5, -0.7);
\node[section label, anchor=west] at (-1.2, -1.2) {Orientabilidad};
\steerablefiltercolor{0.5}{-2.7}{35}{0.65}{gdlorange!70}
\node[below, text=gdlorange!70!black, font=\small] at (0.5, -3.6) {$\psi_\theta$};
\node[font=\Large, text=gdlgray!70!black] at (1.8, -2.7) {$=$};
\node[font=\normalsize, text=gdlred!70!black] at (2.7, -2.7) {$\cos\theta$};
\node[font=\large, text=gdlgray!70!black] at (3.4, -2.7) {$\times$};
\steerablefiltercolor{4.3}{-2.7}{0}{0.55}{gdlred!70}
\node[below, text=gdlred!70!black, font=\small] at (4.3, -3.45) {$\psi_0$};
\node[font=\Large, text=gdlgray!70!black] at (5.4, -2.7) {$+$};
\node[font=\normalsize, text=gdlgreen!70!black] at (6.3, -2.7) {$\sin\theta$};
\node[font=\large, text=gdlgray!70!black] at (7.0, -2.7) {$\times$};
\steerablefiltercolor{7.9}{-2.7}{90}{0.55}{gdlgreen!70}
\node[below, text=gdlgreen!70!black, font=\small] at (7.9, -3.45) {$\psi_{90}$};
\fill[gdlblue!8, rounded corners=4pt] (0.2, -4.4) rectangle (10.3, -5.4);
\draw[gdlblue!40, rounded corners=4pt, thick] (0.2, -4.4) rectangle (10.3, -5.4);
\node[font=\normalsize, text=gdlblue!70!black, anchor=west] at (0.6, -4.9)
    {Orientabilidad:\quad $\displaystyle \rho(r_\theta)\psi = \sum_{k} a_k(\theta)\,\psi_k$};
\node[font=\scriptsize, text=gdlgray!70!black, anchor=west] at (7.0, -4.9)
    {$\{\psi_k\}$: base de filtros};
\end{tikzpicture}
\caption[Filtros orientables]{Filtros orientables: cualquier rotaci\'on del
filtro base $\psi_0$ es una combinaci\'on lineal fija de una base finita de
filtros. Esta propiedad proporciona equivarianza a rotaciones continuas con un
n\'umero finito de par\'ametros.}
\label{fig:groups:steerable}
\end{figure}

En la pr\'actica una convoluci\'on orientable almacena coeficientes aprendibles
que combinan una base fija de arm\'onicos circulares (una envolvente radial
gaussiana por $\cos k\theta$ y $\sin k\theta$ para $k = 0, 1, \dots, F$) en el
n\'ucleo de convoluci\'on. Las no linealidades equivariantes preservan la
estructura irreducible: los canales triviales ($k=0$) usan activaciones puntuales
ordinarias, mientras que los canales no triviales usan una no linealidad de norma
que modula la magnitud de la caracter\'istica preservando su fase. Las redes
orientables que siguen el marco de \citet{weiler2019general} y el aprendizaje
temprano de filtros orientables de \citet{weiler2018learningsteerable} superan a
las discretizaciones finitas en benchmarks rotados usando menos par\'ametros,
porque la base arm\'onica representa la simetr\'ia rotacional continua de forma
exacta y no mediante copias discretas.

\subsection{Representaciones irreducibles y campos geom\'etricos}
\label{sec:groups:irreps-fields}

Las redes orientables operan no sobre canales escalares sino sobre \emph{campos
geom\'etricos} que transforman bajo las representaciones irreducibles de
$\SO(2)$.

\begin{definition}[Campo geom\'etrico]
\label{def:groups:geometric-field}
Un campo geom\'etrico de tipo $\rho$ es una funci\'on $f: \R^2 \to V$ donde $V$
porta una representaci\'on $\rho: \SO(2) \to \GL(V)$ tal que
\begin{equation}
  (g \cdot f)(x) = \rho(g)\, f(g^{-1} x) \qquad \forall g \in \SO(2).
\end{equation}
\end{definition}

\begin{definition}[Representaciones irreducibles de $\SO(2)$]
\label{def:groups:so2-irreps}
Las representaciones irreducibles de $\SO(2)$ est\'an indexadas por enteros $k \in
\Z$: $\rho_k(R_\theta) = e^{ik\theta}$ para $k \neq 0$, y la representaci\'on
trivial $\rho_0(R_\theta) = 1$. Sobre los reales, las frecuencias $\pm k$ se
combinan en la representaci\'on bidimensional
\begin{equation}
  \rho_k^{\text{real}}(R_\theta) = \begin{pmatrix}
    \cos(k\theta) & -\sin(k\theta) \\
    \sin(k\theta) & \cos(k\theta) \end{pmatrix}.
\end{equation}
\end{definition}

Las redes orientables destacan all\'i donde la orientaci\'on exacta es
informativa: an\'alisis de fibras en imagen m\'edica, detecci\'on de crestas en
huellas dactilares, y detecci\'on de grietas y defectos en materiales. Forman el
puente natural hacia el grupo de rotaciones tridimensionales $\SO(3)$, al que
nos volvemos a continuaci\'on.

\section{CNN esf\'ericas y equivarianza a $\SO(3)$}
\label{sec:groups:spherical}

Muchas se\~nales del mundo real residen en la esfera o portan simetr\'ia
rotacional tridimensional completa: las im\'agenes omnidireccionales de
$360^\circ$ est\'an definidas sobre $S^2$; los campos clim\'aticos globales
residen en la superficie (topol\'ogicamente esf\'erica) de la Tierra; el fondo
c\'osmico de microondas es una se\~nal sobre la esfera celeste. Proyectar tales
datos a una rejilla rectangular introduce distorsi\'on de \'area cerca de los
polos, discontinuidades artificiales en las costuras y una p\'erdida de
equivarianza, ya que las rotaciones esf\'ericas se convierten en transformaciones
complicadas en el dominio proyectado. La alternativa fundamentada es construir
convoluciones equivariantes al grupo de rotaciones $\SO(3)$.

\begin{definition}[El grupo $\SO(3)$]
\label{def:groups:so3}
$\SO(3) = \{R \in \R^{3 \times 3} : R^\top R = I,\ \det R = 1\}$ es el grupo de
las rotaciones de $\R^3$ que preservan la orientaci\'on. Cada elemento puede
parametrizarse mediante \'angulos de Euler $(\alpha, \beta, \gamma)$ o mediante un
eje de rotaci\'on $\mathbf{n} \in S^2$ y un \'angulo $\theta \in [0, 2\pi)$.
\end{definition}

\begin{proposition}[Propiedades de $\SO(3)$]
\label{prop:groups:so3-properties}
El grupo $\SO(3)$ es \emph{no abeliano}, \emph{compacto} (topol\'ogicamente el
espacio proyectivo real $\R P^3$) y \emph{triparam\'etrico}.
\end{proposition}

\begin{proof}
\emph{No abeliano:} sean $R_x, R_z$ rotaciones de $90^\circ$ en torno a los ejes
$x$ y $z$; entonces $R_x R_z (1,0,0)^\top = (0,0,1)^\top$ pero $R_z R_x
(1,0,0)^\top = (0,1,0)^\top$, de modo que $R_x R_z \neq R_z R_x$.
\emph{Compacto:} $\SO(3)$ es cerrado en $\R^9$ y acotado ($\norm{R}_F^2 =
\operatorname{tr}(R^\top R) = 3$), y por tanto compacto; el homeomorfismo con
$\R P^3$ proviene de la parametrizaci\'on eje--\'angulo, que identifica $(\pi,
\mathbf{n})$ con $(\pi, -\mathbf{n})$. \emph{Triparam\'etrico:} la condici\'on
$R^\top R = I$ impone seis restricciones independientes sobre las nueve entradas,
de modo que $\dim \SO(3) = 9 - 6 = 3$.
\end{proof}

\subsection{Arm\'onicos esf\'ericos}
\label{sec:groups:spherical-harmonics}

El an\'alisis arm\'onico sobre $S^2$ proporciona la base para la convoluci\'on
equivariante a $\SO(3)$.

\begin{definition}[Arm\'onicos esf\'ericos]
\label{def:groups:spherical-harmonics}
Los arm\'onicos esf\'ericos $Y_\ell^m(\theta, \phi)$ forman una base ortonormal
completa de $L^2(S^2)$:
\begin{equation}
  Y_\ell^m(\theta, \phi) = \sqrt{\frac{2\ell+1}{4\pi}\,
    \frac{(\ell - |m|)!}{(\ell + |m|)!}}\,
    P_\ell^{|m|}(\cos\theta)\, e^{im\phi},
\end{equation}
donde $P_\ell^m$ son los polinomios asociados de Legendre, $\ell \geq 0$ es el
grado, y $|m| \leq \ell$ el orden.
\end{definition}

Toda $f \in L^2(S^2)$ se desarrolla como $f(\theta, \phi) = \sum_{\ell=0}^{\infty}
\sum_{m=-\ell}^{\ell} f_\ell^m\, Y_\ell^m(\theta, \phi)$ con coeficientes
$f_\ell^m = \int_{S^2} f(\mathbf{x})\, \overline{Y_\ell^m(\mathbf{x})}\,
d\mathbf{x}$. \Cref{fig:groups:spherical-harmonics} muestra los arm\'onicos de
grado bajo. Su propiedad definitoria es su comportamiento bajo rotaci\'on.

\begin{theorem}[Transformaci\'on de los arm\'onicos esf\'ericos]
\label{thm:groups:sh-transformation}
Bajo una rotaci\'on $R \in \SO(3)$,
\begin{equation}
  Y_\ell^m(R^{-1}\mathbf{x}) = \sum_{m'=-\ell}^{\ell} D^{(\ell)}_{m'm}(R)\,
    Y_\ell^{m'}(\mathbf{x}),
\end{equation}
donde $D^{(\ell)}(R)$ es la matriz de Wigner irreducible de tama\~no
$(2\ell+1)\times(2\ell+1)$ de $\SO(3)$.
\end{theorem}

\begin{proof}
Los arm\'onicos $\{Y_\ell^m\}_{m=-\ell}^{\ell}$ generan el espacio
$\mathcal{H}_\ell$ de polinomios arm\'onicos homog\'eneos de grado $\ell$
restringidos a $S^2$. Este espacio es $\SO(3)$-invariante, porque la composici\'on
con una rotaci\'on preserva tanto la homogeneidad como la armonicidad (el
laplaciano conmuta con las rotaciones). Como $\mathcal{H}_\ell$ tiene dimensi\'on
$2\ell+1$ y es invariante, $Y_\ell^m \circ R^{-1}$ es una combinaci\'on lineal de
la base, con coeficientes $D^{(\ell)}_{m'm}(R) = \int_{S^2} Y_\ell^m(R^{-1}
\mathbf{x})\, \overline{Y_\ell^{m'}(\mathbf{x})} \, d\mathbf{x}$. La
irreducibilidad de $D^{(\ell)}$ se sigue de que $\mathcal{H}_\ell$ no tiene
subespacio $\SO(3)$-invariante propio.
\end{proof}

\begin{figure}[h]
\centering
\begin{tikzpicture}[every node/.style={font=\small}]
  \def\sphR{1}
  \def\colsep{2.8}
  \def\rowsep{3.2}
  \node[font=\large\bfseries, text=gdlblue!70!black]
    at ({2*\colsep}, 2.4) {Arm\'onicos esf\'ericos $Y_\ell^m(\theta,\varphi)$};
  \foreach \m/\col in {-2/0, -1/1, 0/2, 1/3, 2/4} {
    \node[font=\footnotesize, text=gdlpurple!70!black] at ({\col*\colsep}, 1.5) {$m = \m$};
  }
  \node[font=\footnotesize, text=gdlorange!70!black] at (-1.6, 0)          {$\ell = 0$};
  \node[font=\footnotesize, text=gdlorange!70!black] at (-1.6, -\rowsep)   {$\ell = 1$};
  \node[font=\footnotesize, text=gdlorange!70!black] at (-1.6, -2*\rowsep) {$\ell = 2$};
  % l=0,m=0
  \begin{scope}[shift={(2*\colsep, 0)}]
    \fill[gdlblue!20] (0,0) circle (\sphR);
    \draw[gdlgray!60, thick] (0,0) circle (\sphR);
    \draw[gdlgray!40, dashed] ({\sphR},0) arc[start angle=0, end angle=180, x radius=\sphR, y radius=0.3*\sphR];
    \draw[gdlgray!50] ({-\sphR},0) arc[start angle=180, end angle=360, x radius=\sphR, y radius=0.3*\sphR];
    \node[font=\scriptsize, text=gdlblue!70!black] at (0, -\sphR - 0.4) {$Y_0^0$};
  \end{scope}
  % l=1,m=-1
  \begin{scope}[shift={(1*\colsep, -\rowsep)}]
    \fill[gdlblue!30] (90:\sphR) arc (90:270:\sphR) -- cycle;
    \fill[gdlred!30]  (90:\sphR) arc (90:-90:\sphR) -- cycle;
    \draw[gdlgray!60, thick] (0,0) circle (\sphR);
    \draw[gdlgray!40, dashed] (0,\sphR) -- (0,-\sphR);
    \draw[gdlgray!40, dashed] ({\sphR},0) arc[start angle=0, end angle=180, x radius=\sphR, y radius=0.3*\sphR];
    \draw[gdlgray!50] ({-\sphR},0) arc[start angle=180, end angle=360, x radius=\sphR, y radius=0.3*\sphR];
    \node[font=\scriptsize, text=gdlblue!70!black] at (-0.5, 0.5) {$+$};
    \node[font=\scriptsize, text=gdlred!70!black]  at ( 0.5, 0.5) {$-$};
    \node[font=\scriptsize, text=gdlpurple!70!black] at (0, -\sphR - 0.4) {$Y_1^{-1}$};
  \end{scope}
  % l=1,m=0
  \begin{scope}[shift={(2*\colsep, -\rowsep)}]
    \fill[gdlblue!30] (-\sphR,0) arc (180:0:\sphR) -- cycle;
    \fill[gdlred!30]  (-\sphR,0) arc (180:360:\sphR) -- cycle;
    \draw[gdlgray!60, thick] (0,0) circle (\sphR);
    \draw[gdlgray!50] (-\sphR,0) -- (\sphR,0);
    \draw[gdlgray!40, dashed] ({\sphR},0) arc[start angle=0, end angle=180, x radius=\sphR, y radius=0.3*\sphR];
    \draw[gdlgray!50] ({-\sphR},0) arc[start angle=180, end angle=360, x radius=\sphR, y radius=0.3*\sphR];
    \node[font=\scriptsize, text=gdlblue!70!black] at (0,  0.5) {$+$};
    \node[font=\scriptsize, text=gdlred!70!black]  at (0, -0.5) {$-$};
    \node[font=\scriptsize, text=gdlpurple!70!black] at (0, -\sphR - 0.4) {$Y_1^{0}$};
  \end{scope}
  % l=1,m=1
  \begin{scope}[shift={(3*\colsep, -\rowsep)}]
    \fill[gdlred!30]  (90:\sphR) arc (90:270:\sphR) -- cycle;
    \fill[gdlblue!30] (90:\sphR) arc (90:-90:\sphR) -- cycle;
    \draw[gdlgray!60, thick] (0,0) circle (\sphR);
    \draw[gdlgray!40, dashed] (0,\sphR) -- (0,-\sphR);
    \draw[gdlgray!40, dashed] ({\sphR},0) arc[start angle=0, end angle=180, x radius=\sphR, y radius=0.3*\sphR];
    \draw[gdlgray!50] ({-\sphR},0) arc[start angle=180, end angle=360, x radius=\sphR, y radius=0.3*\sphR];
    \node[font=\scriptsize, text=gdlred!70!black]  at (-0.5, 0.5) {$-$};
    \node[font=\scriptsize, text=gdlblue!70!black] at ( 0.5, 0.5) {$+$};
    \node[font=\scriptsize, text=gdlpurple!70!black] at (0, -\sphR - 0.4) {$Y_1^{1}$};
  \end{scope}
  % l=2,m=-2
  \begin{scope}[shift={(0*\colsep, -2*\rowsep)}]
    \clip (0,0) circle (\sphR);
    \fill[gdlblue!30] (0,0) -- (90:\sphR)  arc (90:0:\sphR)    -- cycle;
    \fill[gdlred!30]  (0,0) -- (0:\sphR)   arc (0:-90:\sphR)   -- cycle;
    \fill[gdlblue!30] (0,0) -- (-90:\sphR) arc (-90:-180:\sphR) -- cycle;
    \fill[gdlred!30]  (0,0) -- (180:\sphR) arc (180:90:\sphR)  -- cycle;
    \draw[gdlgray!40, dashed] (-\sphR,0) -- (\sphR,0);
    \draw[gdlgray!40, dashed] (0,-\sphR) -- (0,\sphR);
    \draw[gdlgray!60, thick] (0,0) circle (\sphR);
    \node[font=\scriptsize, text=gdlpurple!70!black] at (0, -\sphR - 0.4) {$Y_2^{-2}$};
  \end{scope}
  % l=2,m=-1
  \begin{scope}[shift={(1*\colsep, -2*\rowsep)}]
    \clip (0,0) circle (\sphR);
    \fill[gdlblue!30] (0,0) -- (45:\sphR)  arc (45:135:\sphR)  -- cycle;
    \fill[gdlred!30]  (0,0) -- (135:\sphR) arc (135:225:\sphR) -- cycle;
    \fill[gdlblue!30] (0,0) -- (225:\sphR) arc (225:315:\sphR) -- cycle;
    \fill[gdlred!30]  (0,0) -- (315:\sphR) arc (315:405:\sphR) -- cycle;
    \draw[gdlgray!60, thick] (0,0) circle (\sphR);
    \node[font=\scriptsize, text=gdlpurple!70!black] at (0, -\sphR - 0.4) {$Y_2^{-1}$};
  \end{scope}
  % l=2,m=0
  \begin{scope}[shift={(2*\colsep, -2*\rowsep)}]
    \clip (0,0) circle (\sphR);
    \fill[gdlblue!30] (-\sphR, {0.33*\sphR}) rectangle (\sphR, \sphR);
    \fill[gdlred!30]  (-\sphR, {-0.33*\sphR}) rectangle (\sphR, {0.33*\sphR});
    \fill[gdlblue!30] (-\sphR, -\sphR) rectangle (\sphR, {-0.33*\sphR});
    \draw[gdlgray!60, thick] (0,0) circle (\sphR);
    \node[font=\scriptsize, text=gdlpurple!70!black] at (0, -\sphR - 0.4) {$Y_2^{0}$};
  \end{scope}
  % l=2,m=1
  \begin{scope}[shift={(3*\colsep, -2*\rowsep)}]
    \clip (0,0) circle (\sphR);
    \fill[gdlred!30]  (0,0) -- (45:\sphR)  arc (45:135:\sphR)  -- cycle;
    \fill[gdlblue!30] (0,0) -- (135:\sphR) arc (135:225:\sphR) -- cycle;
    \fill[gdlred!30]  (0,0) -- (225:\sphR) arc (225:315:\sphR) -- cycle;
    \fill[gdlblue!30] (0,0) -- (315:\sphR) arc (315:405:\sphR) -- cycle;
    \draw[gdlgray!60, thick] (0,0) circle (\sphR);
    \node[font=\scriptsize, text=gdlpurple!70!black] at (0, -\sphR - 0.4) {$Y_2^{1}$};
  \end{scope}
  % l=2,m=2
  \begin{scope}[shift={(4*\colsep, -2*\rowsep)}]
    \clip (0,0) circle (\sphR);
    \fill[gdlred!30]  (0,0) -- (90:\sphR)  arc (90:0:\sphR)    -- cycle;
    \fill[gdlblue!30] (0,0) -- (0:\sphR)   arc (0:-90:\sphR)   -- cycle;
    \fill[gdlred!30]  (0,0) -- (-90:\sphR) arc (-90:-180:\sphR) -- cycle;
    \fill[gdlblue!30] (0,0) -- (180:\sphR) arc (180:90:\sphR)  -- cycle;
    \draw[gdlgray!60, thick] (0,0) circle (\sphR);
    \node[font=\scriptsize, text=gdlpurple!70!black] at (0, -\sphR - 0.4) {$Y_2^{2}$};
  \end{scope}
  \begin{scope}[shift={(0, {-2*\rowsep - 2.3})}]
    \fill[gdlblue!30] ({0.8*\colsep}, 0) rectangle ++(0.4, 0.4);
    \draw[gdlgray!60] ({0.8*\colsep}, 0) rectangle ++(0.4, 0.4);
    \node[font=\scriptsize, anchor=west, text=gdlblue!70!black]
      at ({0.8*\colsep + 0.55}, 0.2) {Regi\'on positiva ($Y_\ell^m > 0$)};
    \fill[gdlred!30] ({2.5*\colsep}, 0) rectangle ++(0.4, 0.4);
    \draw[gdlgray!60] ({2.5*\colsep}, 0) rectangle ++(0.4, 0.4);
    \node[font=\scriptsize, anchor=west, text=gdlred!70!black]
      at ({2.5*\colsep + 0.55}, 0.2) {Regi\'on negativa ($Y_\ell^m < 0$)};
  \end{scope}
  \node[font=\scriptsize, text=gdlgray!80!black, text width=9cm, align=center]
    at ({2*\colsep}, {-2*\rowsep - 3.3})
    {$\ell$ = grado (frecuencia),\quad $m$ = orden (orientaci\'on)};
\end{tikzpicture}
\caption[Arm\'onicos esf\'ericos]{Arm\'onicos esf\'ericos $Y_\ell^m(\theta,
\varphi)$ para grados $\ell = 0, 1, 2$. El grado $\ell$ controla la frecuencia
angular (n\'umero de cambios de signo) y el orden $m$ la orientaci\'on del
patr\'on nodal. Estas funciones forman la base de las representaciones
irreducibles $(2\ell+1)$-dimensionales de $\SO(3)$.}
\label{fig:groups:spherical-harmonics}
\end{figure}

\subsection{Convoluci\'on esf\'erica}
\label{sec:groups:spherical-convolution}

Las redes convolucionales esf\'ericas, introducidas por
\citet{cohen2018sphericalcnns}, procesan funciones sobre $S^2 = \SO(3)/\SO(2)$
directamente a trav\'es de la estructura de grupo de $\SO(3)$.

\begin{definition}[Convoluci\'on en $\SO(3)$]
\label{def:groups:so3-convolution}
Para $f: \SO(3) \to \R^{C_{\text{in}}}$ y n\'ucleo $\psi: \SO(3) \to
\R^{C_{\text{out}} \times C_{\text{in}}}$,
\begin{equation}
  (f \star \psi)(g) = \int_{\SO(3)} f(h)\, \psi(g^{-1} h)\, dh,
\end{equation}
donde $dh$ es la medida de Haar sobre $\SO(3)$.
\end{definition}

\begin{theorem}[Equivarianza de la convoluci\'on esf\'erica]
\label{thm:groups:spherical-conv-equivariance}
La convoluci\'on en $\SO(3)$ es equivariante: si $(g \cdot f)(x) = f(g^{-1} x)$
entonces $g \cdot (f \star \psi) = (g \cdot f) \star \psi$.
\end{theorem}

\begin{proof}
Calculamos, para $g \in \SO(3)$,
\begin{align}
  [(g \cdot f) \star \psi](x)
  &= \int_{\SO(3)} (g \cdot f)(h)\, \psi(x^{-1} h)\, dh
   = \int_{\SO(3)} f(g^{-1} h)\, \psi(x^{-1} h)\, dh.
\end{align}
Sustituyendo $h' = g^{-1} h$ (de modo que $h = g h'$) y usando la invarianza por
la izquierda de la medida de Haar ($dh = dh'$),
\begin{align}
  &= \int_{\SO(3)} f(h')\, \psi(x^{-1} g h')\, dh'
   = \int_{\SO(3)} f(h')\, \psi\big((g^{-1} x)^{-1} h'\big)\, dh'
   = (f \star \psi)(g^{-1} x),
\end{align}
que es exactamente $[g \cdot (f \star \psi)](x)$.
\end{proof}

Las implementaciones eficientes usan una transformada r\'apida de Fourier
esf\'erica: el teorema de convoluci\'on se extiende a la esfera, reduciendo la
convoluci\'on a una multiplicaci\'on de matrices dentro de cada bloque de grado
$\ell$, $\widehat{(f \star \psi)}_\ell^m = \sum_{m'} \hat f_\ell^{m'}
\hat\psi_\ell^{m'm}$. La complejidad est\'a gobernada por el grado m\'aximo $L$:
$O(L^2)$ coeficientes por ubicaci\'on esf\'erica y $O(L^3)$ para la propia
transformada. La construcci\'on se extiende a datos volum\'etricos combinando el
an\'alisis en arm\'onicos esf\'ericos con perfiles radiales, como en las redes
orientables 3D de \citet{weiler2018steerable}, y a rejillas esf\'ericas no
estructuradas \citep{cohen2021spherical}. La progresi\'on desde las redes
convolucionales est\'andar a las redes convolucionales de grupo discreto, a las
redes orientables $\SO(2)$, a las redes esf\'ericas $\SO(3)$ traza un aumento
constante de la conciencia geom\'etrica, ajustando cada nivel la arquitectura m\'as
estrechamente a la simetr\'ia intr\'inseca de los datos.

\section{Grupos no compactos: equivarianza a $\Eucl(3)$ y $\SE(3)$}
\label{sec:groups:noncompact}

Las nubes de puntos, las mol\'eculas y los sistemas de part\'iculas portan
simetr\'ias que incluyen \emph{traslaciones} adem\'as de rotaciones. Los grupos
relevantes son el grupo eucl\'ideo $\Eucl(n)$ y el grupo eucl\'ideo especial
$\SE(n)$, que son \emph{no compactos}: el subgrupo de traslaciones $\R^n$ es no
acotado, de modo que la medida de Haar es infinita y la integraci\'on directa
sobre el grupo deja de ser factible. Las arquitecturas equivariantes deben, por
tanto, construirse de otra manera.

\begin{definition}[Grupos eucl\'ideo y eucl\'ideo especial]
\label{def:groups:euclidean}
El \emph{grupo eucl\'ideo} $\Eucl(n) = \R^n \rtimes \OO(n)$ es el grupo de todas
las isometr\'ias de $\R^n$. Cada elemento es un par $(R, \mathbf{t})$ con $R \in
\OO(n)$ y $\mathbf{t} \in \R^n$, que act\'ua mediante $g \cdot \mathbf{x} =
R\mathbf{x} + \mathbf{t}$, con composici\'on $(R_1, \mathbf{t}_1)\circ(R_2,
\mathbf{t}_2) = (R_1 R_2,\, R_1 \mathbf{t}_2 + \mathbf{t}_1)$. El \emph{grupo
eucl\'ideo especial} $\SE(n) = \R^n \rtimes \SO(n)$ es el subgrupo que preserva
la orientaci\'on.
\end{definition}

\begin{proposition}[Propiedades de $\Eucl(n)$]
\label{prop:groups:en-properties}
$\Eucl(n)$ es \emph{no compacto} (su subgrupo de traslaciones $\R^n$ es no
acotado), un \emph{grupo de Lie} de dimensi\'on $n + \tfrac{n(n-1)}{2}$, y un
\emph{producto semidirecto} en el que $\R^n$ es un subgrupo normal.
\end{proposition}

\begin{proof}
\emph{No compacto:} las traslaciones puras $\{(I, \mathbf{t})\} \cong \R^n$ son
no acotadas. \emph{Grupo de Lie:} $\Eucl(n)$ se sumerge como subgrupo cerrado de
$\GL(n+1, \R)$ v\'ia $(R, \mathbf{t}) \mapsto \left(\begin{smallmatrix} R &
\mathbf{t} \\ 0 & 1 \end{smallmatrix}\right)$, de dimensi\'on $n + \dim \OO(n) =
n + \tfrac{n(n-1)}{2}$. \emph{Producto semidirecto:} al conjugar una traslaci\'on
pura, $(R, \mathbf{s})(I, \mathbf{t})(R, \mathbf{s})^{-1} = (I, R\mathbf{t})$,
sigue siendo una traslaci\'on, de modo que $\R^n$ es normal y la ley de
composici\'on es exactamente la de $\R^n \rtimes \OO(n)$.
\end{proof}

El hecho estructural crucial para el dise\~no de arquitecturas es que una
representaci\'on de dimensi\'on finita en la que las traslaciones act\'uan
trivialmente se factoriza a trav\'es de la parte de rotaci\'on.

\begin{proposition}[Descomposici\'on de la representaci\'on]
\label{prop:groups:rep-decomposition}
En cualquier representaci\'on de dimensi\'on finita de $\Eucl(n) = \R^n \rtimes
\OO(n)$ en la que las traslaciones act\'uan trivialmente, la representaci\'on se
factoriza a trav\'es del cociente $\Eucl(n)/\R^n \cong \OO(n)$: $\rho(R,
\mathbf{t}) = \rho_{\OO(n)}(R)$. El enunciado an\'alogo vale para $\SE(n)$ a
trav\'es de $\SO(n)$ \citep{DBLP:journals/corr/abs-2104-13478}.
\end{proposition}

\begin{proof}
Sea $\rho: \Eucl(n) \to \GL(V)$ que satisface $\rho(I, \mathbf{t}) =
\operatorname{Id}$ para todo $\mathbf{t}$. Como $(R, \mathbf{t}) = (R,
\mathbf{0})\cdot(I, R^{-1}\mathbf{t})$, tenemos $\rho(R, \mathbf{t}) = \rho(R,
\mathbf{0})$. Definiendo $\rho_{\OO(n)}(R) := \rho(R, \mathbf{0})$ se obtiene un
homomorfismo, porque $\rho_{\OO(n)}(R_1 R_2) = \rho(R_1 R_2, \mathbf{0}) =
\rho(R_1, \mathbf{0}) \rho(R_2, \mathbf{0})$. As\'i, $\rho$ se factoriza a
trav\'es de la proyecci\'on can\'onica $\pi: \Eucl(n) \to \Eucl(n)/\R^n \cong
\OO(n)$.
\end{proof}

Esta factorizaci\'on motiva arquitecturas que separan la informaci\'on
\emph{invariante} (caracter\'isticas escalares) de la informaci\'on
\emph{equivariante} (coordenadas vectoriales), y da lugar a dos familias de
modelos: redes de grafos construidas sobre escalares invariantes, y redes de
campos tensoriales construidas sobre representaciones completas de $\SO(3)$.

\subsection{Mensajes invariantes: redes de grafos equivariantes a \texorpdfstring{$\Eucl(n)$}{E(n)}}
\label{sec:groups:egnn}

Las redes neuronales de grafos equivariantes a $\Eucl(n)$ de
\citet{satorras2021enequivariant} (EGNN) construyen mensajes a partir de
cantidades $\Eucl(n)$-invariantes mientras actualizan las coordenadas mediante
diferencias vectoriales que preservan la equivarianza. Cada nodo porta
caracter\'isticas escalares $h_i \in \R^d$ (invariantes bajo $\Eucl(n)$) y
coordenadas $\mathbf{x}_i \in \R^n$ (equivariantes bajo $\Eucl(n)$), que se
actualizan en tres etapas.

\paragraph{Etapa 1: mensajes invariantes.}
Los mensajes dependen solo de cantidades $\Eucl(n)$-invariantes:
\begin{equation}
  \mathbf{m}_{ij} = \phi_e\big(h_i, h_j, \norm{\mathbf{x}_i - \mathbf{x}_j}^2,
    a_{ij}\big),
  \label{eq:groups:egnn-message}
\end{equation}
donde $h_i, h_j$ son caracter\'isticas escalares, $\norm{\mathbf{x}_i -
\mathbf{x}_j}^2$ es la distancia eucl\'idea al cuadrado, $a_{ij}$ son atributos
opcionales de arista, y $\phi_e$ es un perceptr\'on multicapa. La distancia al
cuadrado es $\Eucl(n)$-invariante: para cualquier $(R, \mathbf{t})$,
\begin{equation}
  \norm{(R\mathbf{x}_i + \mathbf{t}) - (R\mathbf{x}_j + \mathbf{t})}^2
  = \norm{R(\mathbf{x}_i - \mathbf{x}_j)}^2
  = \norm{\mathbf{x}_i - \mathbf{x}_j}^2,
\end{equation}
puesto que $R$ es ortogonal y las traslaciones se cancelan.

\paragraph{Etapa 2: actualizaci\'on equivariante de coordenadas.}
Las coordenadas se actualizan mediante diferencias vectoriales ponderadas:
\begin{equation}
  \mathbf{x}_i' = \mathbf{x}_i + C \sum_{j \in \mathcal{N}(i)}
    (\mathbf{x}_i - \mathbf{x}_j)\, \phi_x(\mathbf{m}_{ij}),
  \label{eq:groups:egnn-coord}
\end{equation}
donde $\phi_x$ es un perceptr\'on multicapa de valores escalares, $C$ una
constante de normalizaci\'on, y $\mathcal{N}(i)$ la vecindad del nodo $i$.

\begin{proposition}[Equivarianza de la actualizaci\'on de coordenadas]
\label{prop:groups:egnn-equivariance}
La actualizaci\'on \Cref{eq:groups:egnn-coord} es $\Eucl(n)$-equivariante:
aplicar $(R, \mathbf{t})$ a todas las coordenadas produce $\mathbf{x}_i' \mapsto
R\mathbf{x}_i' + \mathbf{t}$.
\end{proposition}

\begin{proof}
Con coordenadas transformadas $\tilde{\mathbf{x}}_i = R\mathbf{x}_i +
\mathbf{t}$, los mensajes $\mathbf{m}_{ij}$ no cambian (est\'an construidos a
partir de invariantes), de modo que
\begin{align}
  \tilde{\mathbf{x}}_i'
  &= R\mathbf{x}_i + \mathbf{t} + C \sum_{j} \big((R\mathbf{x}_i + \mathbf{t})
     - (R\mathbf{x}_j + \mathbf{t})\big)\, \phi_x(\mathbf{m}_{ij}) \\
  &= R\mathbf{x}_i + \mathbf{t} + C \sum_{j} R(\mathbf{x}_i - \mathbf{x}_j)\,
     \phi_x(\mathbf{m}_{ij})
   = R\Big(\mathbf{x}_i + C \sum_{j} (\mathbf{x}_i - \mathbf{x}_j)\,
     \phi_x(\mathbf{m}_{ij})\Big) + \mathbf{t}
   = R\mathbf{x}_i' + \mathbf{t}. \qedhere
\end{align}
\end{proof}

\paragraph{Etapa 3: actualizaci\'on invariante de caracter\'isticas.}
Las caracter\'isticas escalares se actualizan a partir de los mensajes agregados,
$\mathbf{m}_i = \sum_{j \in \mathcal{N}(i)} \mathbf{m}_{ij}$ y $h_i' = \phi_h(h_i,
\mathbf{m}_i)$, lo cual es trivialmente invariante porque opera \'unicamente sobre
cantidades invariantes. La arquitectura garantiza as\'i $h_i'(g \cdot \mathbf{X})
= h_i'(\mathbf{X})$ (invarianza) y $\mathbf{x}_i'(g \cdot \mathbf{X}) = g \cdot
\mathbf{x}_i'(\mathbf{X})$ (equivarianza) para todo $g = (R, \mathbf{t})$.

\begin{example}[Paso de mensajes EGNN para una mol\'ecula peque\~na]
\label{ex:groups:egnn-molecular}
Consid\'erense tres \'atomos en $\mathbf{x}_1 = (0,0,0)$, $\mathbf{x}_2 =
(1.5,0,0)$, $\mathbf{x}_3 = (0.75, 1.3, 0)$ con caracter\'isticas escalares $h_1
= [6.0]$ (carbono), $h_2 = [1.0]$ (hidr\'ogeno), $h_3 = [8.0]$ (ox\'igeno). Las
distancias al cuadrado son $d_{12}^2 = 2.25$, $d_{13}^2 = 0.75^2 + 1.3^2 =
2.2525$, $d_{23}^2 = 2.2525$. Los mensajes son $\mathbf{m}_{12} = \phi_e(h_1,
h_2, d_{12}^2) = \MLP([6.0, 1.0, 2.25])$, y la actualizaci\'on de coordenadas es
$\mathbf{x}_1' = \mathbf{x}_1 + \tfrac{1}{2}[(\mathbf{x}_1 -
\mathbf{x}_2)\phi_x(\mathbf{m}_{12}) + (\mathbf{x}_1 -
\mathbf{x}_3)\phi_x(\mathbf{m}_{13})]$. Rotar la mol\'ecula completa rota las
coordenadas actualizadas de forma id\'entica, por
\Cref{prop:groups:egnn-equivariance}.
\end{example}

Las EGNN son altamente escalables y adecuadas para sistemas grandes (miles de
\'atomos), a costa de la expresividad: usan solo caracter\'isticas escalares y
vectores de posici\'on, de modo que no pueden representar interacciones
direccionales de orden superior. Incorporar invariantes adicionales ---\'angulos
de triplete $\theta_{ijk}$, \'angulos diedros $\phi_{ijkl}$ (que son
$\SE(3)$-invariantes pero cambian de signo bajo reflexi\'on, y as\'i distinguen
la quiralidad), vol\'umenes orientados, o cantidades f\'isicas como los
potenciales de Coulomb y de Lennard-Jones--- mejora el rendimiento sin romper la
simetr\'ia \citep{tholke2022geometric}. Para equivarianza $\Eucl(3)$ completa,
incluyendo reflexiones, se usa $|\phi_{ijkl}|$ en lugar del \'angulo diedro con
signo.

\subsection{Caracter\'isticas tensoriales: Tensor Field Networks y $\SE(3)$-Transformers}
\label{sec:groups:tfn}

Una familia m\'as expresiva representa las caracter\'isticas de nodo como
tensores completos que transforman bajo las representaciones irreducibles de
$\SO(3)$.

\begin{definition}[Representaciones irreducibles de $\SO(3)$]
\label{def:groups:so3-irreps}
Las representaciones irreducibles de $\SO(3)$ est\'an indexadas por enteros no
negativos $\ell = 0, 1, 2, \dots$ (el grado, o momento angular). La
representaci\'on de grado $\ell$ tiene dimensi\'on $2\ell + 1$, con elementos de
matriz, en \'angulos de Euler $(\alpha, \beta, \gamma)$,
\begin{equation}
  D^{(\ell)}_{mm'}(\alpha, \beta, \gamma)
  = e^{-im\alpha}\, d^{(\ell)}_{mm'}(\beta)\, e^{-im'\gamma},
\end{equation}
donde $d^{(\ell)}_{mm'}$ son las matrices peque\~nas de Wigner.
\end{definition}

Los grados bajos tienen un significado geom\'etrico directo: $\ell = 0$ son
escalares (invariantes), $\ell = 1$ son vectores ($\mathbf{v}' = R\mathbf{v}$), y
$\ell = 2$ son tensores sim\'etricos de segundo orden sin traza. Un \emph{tensor
de caracter\'isticas} de grado $\ell$, $\mathbf{f}_i^{(\ell)} \in \R^{2\ell+1}$,
en el nodo $i$ transforma como $(\mathbf{f}_i^{(\ell)})' = D^{(\ell)}(R)\,
\mathbf{f}_i^{(\ell)}$, y cada nodo contiene varias copias (multiplicidades) de
cada tipo, $\mathbf{F}_i = \bigoplus_{\ell=0}^{L_{\max}} \bigoplus_{k=1}^{n_\ell}
\mathbf{f}_{i,k}^{(\ell)}$.

\Cref{fig:groups:so3} recuerda las matrices de rotaci\'on elementales y las
propiedades b\'asicas de $\SO(3)$. Combinar caracter\'isticas de distintos grados
requiere el producto tensorial de Clebsch--Gordan.

\begin{figure}[h]
\centering
\begin{tikzpicture}[>=Stealth, scale=1.1]
\draw[thick, ->] (0,0,0) -- (3,0,0) node[right] {$x$};
\draw[thick, ->] (0,0,0) -- (0,3,0) node[above] {$y$};
\draw[thick, ->] (0,0,0) -- (0,0,3) node[below left] {$z$};
\draw[gdlgray!40, dashed] (0,0) circle (2cm);
\draw[->, very thick, gdlblue!70] (0,0,0) -- (2,0,0);
\fill[gdlblue!70] (2,0,0) circle (3pt);
\node[gdlblue!70, right] at (2.1,0,0) {$\mathbf{v}$};
\draw[->, very thick, gdlred!70] (0,0,0) -- ({2*cos(45)},{2*sin(45)},0);
\fill[gdlred!70] ({2*cos(45)},{2*sin(45)},0) circle (3pt);
\node[gdlred!70, above right] at ({2*cos(45)},{2*sin(45)},0) {$R_z(\theta)\mathbf{v}$};
\draw[thick, gdlorange!70, ->] (1.2,0,0) arc (0:45:1.2cm);
\node[gdlorange!70] at (1, 0.5, 0) {$\theta$};
\node[font=\small, text width=7cm] at (7, 0.8) {
    \textbf{Rotaciones elementales:}\\[0.4em]
    $R_x(\theta) = \begin{psmallmatrix} 1 & 0 & 0 \\ 0 & \cos\theta & -\sin\theta \\ 0 & \sin\theta & \cos\theta \end{psmallmatrix}$,\quad
    $R_y(\theta) = \begin{psmallmatrix} \cos\theta & 0 & \sin\theta \\ 0 & 1 & 0 \\ -\sin\theta & 0 & \cos\theta \end{psmallmatrix}$,\quad
    $R_z(\theta) = \begin{psmallmatrix} \cos\theta & -\sin\theta & 0 \\ \sin\theta & \cos\theta & 0 \\ 0 & 0 & 1 \end{psmallmatrix}$
};
\node[font=\small, text width=6cm] at (7, -2.4) {
    \textbf{Propiedades de $\SO(3)$:}
    \begin{itemize}
        \item Grupo de Lie compacto, dimensi\'on $3$
        \item $\det R = 1$, $R^\top R = I$
        \item No abeliano: $R_x R_y \neq R_y R_x$
        \item Topolog\'ia: $\R P^3$
    \end{itemize}
};
\end{tikzpicture}
\caption[Rotaciones en $\SO(3)$]{Rotaciones en tres dimensiones. Todo $R \in
\SO(3)$ es una rotaci\'on en torno a alg\'un eje $\mathbf{n}$ por un \'angulo
$\theta$, y se descompone en rotaciones elementales en torno a los ejes
coordenados.}
\label{fig:groups:so3}
\end{figure}

\begin{theorem}[Descomposici\'on de Clebsch--Gordan]
\label{thm:groups:clebsch-gordan}
El producto tensorial de dos representaciones irreducibles de $\SO(3)$ se
descompone como una suma directa de irreducibles:
\begin{equation}
  \rho^{(\ell_1)} \otimes \rho^{(\ell_2)}
  = \bigoplus_{\ell = |\ell_1 - \ell_2|}^{\ell_1 + \ell_2} \rho^{(\ell)}.
\end{equation}
Los coeficientes de Clebsch--Gordan $C^{\ell m}_{\ell_1 m_1, \ell_2 m_2}$
implementan la descomposici\'on:
\begin{equation}
  (f^{(\ell_1)} \otimes g^{(\ell_2)})^{(\ell)}_m
  = \sum_{m_1, m_2} C^{\ell m}_{\ell_1 m_1, \ell_2 m_2}\,
    f^{(\ell_1)}_{m_1}\, g^{(\ell_2)}_{m_2}.
\end{equation}
\end{theorem}

La regla de selecci\'on $|\ell_1 - \ell_2| \leq \ell \leq \ell_1 + \ell_2$
controla qu\'e grados aparecen. Por ejemplo, un producto vector $\times$ vector
($\ell = 1 \times \ell = 1$) produce $\ell \in \{0, 1, 2\}$: un escalar (el
producto escalar), un pseudovector (el producto vectorial), y un tensor
sim\'etrico sin traza.

\paragraph{Tensor Field Networks.}
Las Tensor Field Networks de \citet{thomas2018tensorfieldnetworks} proporcionan
la convoluci\'on $\SE(3)$-equivariante fundacional. Una convoluci\'on aplica un
campo tensorial de entrada de grado $\ell_{\text{in}}$ a una salida de grado
$\ell_{\text{out}}$ mediante un n\'ucleo que depende de la posici\'on relativa
$\mathbf{r}_{ij} = \mathbf{x}_i - \mathbf{x}_j$, factorizado en una parte radial
aprendible y una parte angular fija dada por arm\'onicos esf\'ericos:
\begin{equation}
  K^{(\ell_{\text{out}}, \ell_{\text{in}})}_{ij}
  = \sum_{\ell_f} R^{\ell_f}_{\ell_{\text{out}}, \ell_{\text{in}}}(r_{ij})\,
    \mathcal{Y}^{(\ell_f)}(\hat{\mathbf{r}}_{ij}),
  \label{eq:groups:tfn-kernel}
\end{equation}
donde $r_{ij} = \norm{\mathbf{r}_{ij}}$, $\hat{\mathbf{r}}_{ij} =
\mathbf{r}_{ij}/r_{ij}$, los $\mathcal{Y}^{(\ell_f)}$ son arm\'onicos
esf\'ericos, y los $R^{\ell_f}_{\ell_{\text{out}}, \ell_{\text{in}}}$ son
funciones radiales aprendibles (capas gaussianas, funciones de Bessel, o
envolventes polin\'omicas).

\begin{proposition}[Equivarianza del n\'ucleo factorizado]
\label{prop:groups:tfn-kernel-equivariance}
El n\'ucleo \Cref{eq:groups:tfn-kernel} es $\SE(3)$-equivariante: bajo una
rotaci\'on $R \in \SO(3)$,
\begin{equation}
  K^{(\ell_{\text{out}}, \ell_{\text{in}})}_{ij}(R\mathbf{r}_{ij})
  = D^{(\ell_{\text{out}})}(R)\,
    K^{(\ell_{\text{out}}, \ell_{\text{in}})}_{ij}(\mathbf{r}_{ij})\,
    \big(D^{(\ell_{\text{in}})}(R)\big)^{-1}.
\end{equation}
\end{proposition}

\begin{proof}
La parte radial es invariante a rotaciones ($\norm{R\mathbf{r}} =
\norm{\mathbf{r}}$), y los arm\'onicos esf\'ericos transforman como
$Y_\ell^m(R^{-1}\hat{\mathbf{r}}) = \sum_{m'} D^{(\ell)}_{m'm}(R)\,
Y_\ell^{m'}(\hat{\mathbf{r}})$ (\Cref{thm:groups:sh-transformation}). Acoplar los
arm\'onicos mediante los coeficientes de Clebsch--Gordan traduce esta
transformaci\'on en la ley de transformaci\'on tensorial
$D^{(\ell_{\text{out}})}(R)\, K\, (D^{(\ell_{\text{in}})}(R))^{-1}$, como se
requiere para la equivarianza. Las traslaciones no afectan al n\'ucleo, que
depende solo de la posici\'on relativa.
\end{proof}

El acoplamiento entre un tensor de entrada y la parte angular se realiza mediante
el producto de Clebsch--Gordan, con la regla de selecci\'on $|\ell_{\text{in}} -
\ell_f| \leq \ell_{\text{out}} \leq \ell_{\text{in}} + \ell_f$ determinando los
acoplamientos permitidos. Una capa TFN completa suma estos acoplamientos sobre
los grados de entrada, los grados del filtro y los vecinos:
\begin{equation}
  T_i^{(\ell), \text{out}}
  = \sum_{\ell'} \sum_{j \in \mathcal{N}(i)} \sum_{\ell_f}
    W^{\ell_f}_{\ell, \ell'}(r_{ij})
    \big(T_j^{(\ell')} \otimes \mathcal{Y}^{(\ell_f)}(\hat{\mathbf{r}}_{ij})
    \big)^{(\ell)}.
\end{equation}

\paragraph{$\SE(3)$-Transformers.}
Los $\SE(3)$-Transformers de \citet{fuchs2020se3transformers} combinan estas
convoluciones de campos tensoriales con atenci\'on, dando una arquitectura m\'as
expresiva que las EGNN. Los pesos de atenci\'on se calculan \'unicamente a partir
de cantidades invariantes ---las componentes escalares ($\ell = 0$) y la
distancia por pares:
\begin{equation}
  \alpha_{ij} = \softmax_j\Big(\MLP\big(\mathbf{f}_i^{(0)} \oplus
    \mathbf{f}_j^{(0)} \oplus h(r_{ij})\big)\Big),
  \label{eq:groups:se3-attention}
\end{equation}
donde $r_{ij} = \norm{\mathbf{x}_i - \mathbf{x}_j}$, $h(r_{ij})$ es una
incrustaci\'on radial, y $\oplus$ denota concatenaci\'on.

\begin{proposition}[Invarianza de los pesos de atenci\'on]
\label{prop:groups:se3-attention-invariance}
Los pesos $\alpha_{ij}$ de \Cref{eq:groups:se3-attention} son
$\SE(3)$-invariantes.
\end{proposition}

\begin{proof}
Los pesos dependen de las posiciones solo a trav\'es de $\mathbf{x}_i -
\mathbf{x}_j$ y su norma. Bajo $(R, \mathbf{t}) \in \SE(3)$, la diferencia se
convierte en $R(\mathbf{x}_i - \mathbf{x}_j)$ y la norma se preserva por la
ortogonalidad. Como $\alpha_{ij}$ es una funci\'on de $\norm{\mathbf{x}_i -
\mathbf{x}_j}$ y de las caracter\'isticas escalares invariantes, es invariante.
\end{proof}

La actualizaci\'on de caracter\'isticas aplica estos pesos invariantes a mensajes
de valor construidos mediante n\'ucleos equivariantes al estilo TFN,
$\mathbf{f}_i^{(\ell),\text{new}} = \sum_{j} \alpha_{ij}\,
\mathbf{v}_j^{(\ell)}$, donde cada valor $\mathbf{v}_j^{(\ell)}$ transforma como
un tensor de grado $\ell$.

\begin{theorem}[Equivarianza de los $\SE(3)$-Transformers]
\label{thm:groups:se3-transformer-equivariance}
La actualizaci\'on $\mathbf{f}_i^{(\ell),\text{new}} = \sum_{j} \alpha_{ij}\,
\mathbf{v}_j^{(\ell)}$ con pesos invariantes \Cref{eq:groups:se3-attention} es
$\SE(3)$-equivariante: $(\mathbf{f}_i^{(\ell),\text{new}})' = D^{(\ell)}(R)\,
\mathbf{f}_i^{(\ell),\text{new}}$.
\end{theorem}

\begin{proof}
Los pesos $\alpha_{ij}$ son invariantes
(\Cref{prop:groups:se3-attention-invariance}), y los mensajes de valor,
construidos a partir de n\'ucleos TFN
\citep{thomas2018tensorfieldnetworks}, transforman como
$(\mathbf{v}_j^{(\ell)})' = D^{(\ell)}(R)\, \mathbf{v}_j^{(\ell)}$. Por tanto
\begin{align}
  (\mathbf{f}_i^{(\ell),\text{new}})'
  = \sum_{j} \alpha_{ij}\, D^{(\ell)}(R)\, \mathbf{v}_j^{(\ell)}
  = D^{(\ell)}(R) \sum_{j} \alpha_{ij}\, \mathbf{v}_j^{(\ell)}
  = D^{(\ell)}(R)\, \mathbf{f}_i^{(\ell),\text{new}},
\end{align}
usando la invarianza de $\alpha_{ij}$ y la linealidad de $D^{(\ell)}(R)$. Las
traslaciones act\'uan solo sobre las posiciones, no sobre las caracter\'isticas
tensoriales, completando la equivarianza $\SE(3)$.
\end{proof}

Las dos familias ocupan puntos complementarios en el espectro
expresividad--coste, resumido en \Cref{tab:groups:so3-vs-se3}. El grupo compacto
$\SO(3)$ admite una medida de Haar finita normalizable y campos tensoriales puros
sin invariantes intr\'insecos, mientras que el $\SE(3)$ no compacto tiene una
medida de Haar infinita e invariantes naturales (distancias relativas) que
permiten el procesamiento escalar con perceptrones multicapa est\'andar. Las EGNN
son m\'as baratas y escalan a miles de \'atomos; las TFN y los
$\SE(3)$-Transformers son m\'as expresivos pero, con productos de Clebsch--Gordan
de coste $O(L_{\max}^3)$, est\'an limitados a sistemas medianos con $L_{\max} \leq
4$. El principal cuello de botella computacional de los enfoques tensoriales es el
coste de los productos de Clebsch--Gordan, que crece con $L_{\max}$, junto con la
estabilidad num\'erica de las sumas de arm\'onicos esf\'ericos de grado alto.

\begin{table}[h]
\centering
\caption{Arquitecturas equivariantes: $\SO(3)$ compacto frente a $\SE(3)$ no
compacto.}
\label{tab:groups:so3-vs-se3}
\begin{tabular}{@{}lp{4.6cm}p{4.6cm}@{}}
\toprule
\textbf{Propiedad} & \textbf{$\SO(3)$ (compacto)} & \textbf{$\SE(3)$ (no compacto)} \\
\midrule
Estructura de grupo & Solo rotaciones & Rotaciones y traslaciones \\
\addlinespace
Compacidad & Compacto ($\sim \R P^3$) & No compacto (subgrupo $\R^3$ no acotado) \\
\addlinespace
Medida de Haar & Finita, normalizable & Infinita (no normalizable sobre $\R^3$) \\
\addlinespace
Irreducibles & Todas de dimensi\'on finita, $\dim \rho^{(\ell)} = 2\ell+1$ & Finitas (parte $\SO(3)$) e infinitas (parte $\R^3$) \\
\addlinespace
Invariantes b\'asicos & Ninguno (transitivo sobre $S^2$) & Distancias $\norm{\mathbf{x}_i - \mathbf{x}_j}$ \\
\addlinespace
Arquitectura & Campos tensoriales puros & Separaci\'on escalar/vector esencial \\
\addlinespace
Aplicaciones & Datos esf\'ericos, CMB, clima & Mol\'eculas, nubes de puntos, rob\'otica \\
\bottomrule
\end{tabular}
\end{table}

Las redes equivariantes a $\Eucl(n)$ y $\SE(3)$ aportan mejoras consistentes en
eficiencia de datos (a menudo una reducci\'on de un orden de magnitud en los
requisitos de entrenamiento), generalizaci\'on fuera de distribuci\'on, y
transferibilidad a configuraciones no vistas en dominios donde hay simetr\'ias
eucl\'ideas o r\'igidas: sistemas moleculares, nubes de puntos tridimensionales,
y an\'alisis de part\'iculas.

\section{Grupos de permutaciones: aprender sobre conjuntos}
\label{sec:groups:permutations}

El grupo sim\'etrico $\Sym_n$ es la simetr\'ia de las colecciones no ordenadas:
un conjunto no tiene orden intr\'inseco, de modo que una red que lo procesa debe
ser insensible al reetiquetado. A diferencia de las rotaciones o las
traslaciones, las permutaciones son puramente combinatorias.

\begin{definition}[Grupo sim\'etrico]
\label{def:groups:symmetric-group}
El \emph{grupo sim\'etrico} $\Sym_n$ es el grupo de todas las biyecciones
(permutaciones) de $\{1, \dots, n\}$ bajo composici\'on.
\end{definition}

\begin{proposition}[Propiedades b\'asicas de $\Sym_n$]
\label{prop:groups:sn-properties}
$\Sym_n$ tiene orden $n!$, es no abeliano para $n \geq 3$, est\'a generado por
las transposiciones adyacentes $(i, i+1)$, y admite la presentaci\'on de Coxeter
\begin{equation}
  \Sym_n = \langle \sigma_1, \dots, \sigma_{n-1} \mid \sigma_i^2 = e,\,
    (\sigma_i \sigma_{i+1})^3 = e,\, (\sigma_i \sigma_j)^2 = e \text{ para }
    |i-j| > 1 \rangle.
\end{equation}
\end{proposition}

\begin{proof}[Parcial]
Hay $n$ opciones para la imagen de $1$, $n-1$ para $2$, y as\'i sucesivamente,
dando $n!$ permutaciones. Para $n = 3$, $(12)(23) = (123) \neq (132) = (23)(12)$,
de modo que $\Sym_n$ es no abeliano. La generaci\'on por transposiciones y la
presentaci\'on se establecen por construcci\'on expl\'icita; v\'ease
\citet{dummit2004abstract}.
\end{proof}

Para el c\'omputo, las permutaciones se realizan como matrices.

\begin{definition}[Matriz de permutaci\'on]
\label{def:groups:permutation-matrix}
Para $\pi \in \Sym_n$, la matriz de permutaci\'on $P_\pi \in \R^{n \times n}$
tiene $(P_\pi)_{ij} = 1$ si $j = \pi(i)$ y $0$ en caso contrario.
\end{definition}

\begin{proposition}[Propiedades de las matrices de permutaci\'on]
\label{prop:groups:permutation-matrix-properties}
Las matrices de permutaci\'on satisfacen $P_{\pi_1 \pi_2} = P_{\pi_1} P_{\pi_2}$
(un homomorfismo), $P_\pi^{-1} = P_\pi^\top$ (ortogonalidad), y $P_\pi \mathbf{x}$
permuta las componentes de $\mathbf{x}$.
\end{proposition}

\begin{proof}
Escribiendo $(P_\pi)_{ij} = \delta_{i, \pi(j)}$, $(P_{\pi_1} P_{\pi_2})_{ij} =
\sum_k \delta_{i, \pi_1(k)} \delta_{k, \pi_2(j)} = \delta_{i, \pi_1(\pi_2(j))} =
(P_{\pi_1 \pi_2})_{ij}$. A continuaci\'on $(P_\pi^\top)_{ij} = \delta_{j, \pi(i)}
= (P_{\pi^{-1}})_{ij}$, y $P_\pi P_{\pi^{-1}} = P_e = I$ da $P_\pi^{-1} =
P_\pi^\top$. Finalmente $(P_\pi \mathbf{x})_i = \sum_j \delta_{i, \pi(j)} x_j =
x_{\pi^{-1}(i)}$.
\end{proof}

\begin{theorem}[Teorema de Cayley para $\Sym_n$]
\label{thm:groups:cayley}
Todo grupo finito de orden $n$ es isomorfo a un subgrupo de $\Sym_n$. En
particular $\Sym_n$ contiene una copia de todo grupo finito de orden a lo sumo
$n!$.
\end{theorem}

El teorema de Cayley implica que cualquier simetr\'ia discreta finita puede
codificarse como un subgrupo de permutaciones, lo que justifica el estudio
general de las arquitecturas $\Sym_n$-invariantes y $\Sym_n$-equivariantes.

\paragraph{Funciones invariantes y equivariantes sobre conjuntos.}
Una funci\'on $f: \R^{n \times d} \to \R^k$ es \emph{invariante a permutaciones}
si $f(\pi \cdot X) = f(X)$ para todo $\pi \in \Sym_n$, donde $\pi \cdot X$
permuta las filas de $X$; es \emph{equivariante a permutaciones} si $f(\pi \cdot
X) = \pi \cdot f(X)$. Entre los invariantes b\'asicos est\'an la suma $\sum_i
\mathbf{x}_i$, el m\'aximo elemento a elemento, y la media $\tfrac{1}{n}\sum_i
\mathbf{x}_i$.

\paragraph{Teor\'ia de representaciones de $\Sym_n$.}
Las clases de conjugaci\'on de $\Sym_n$ est\'an en biyecci\'on con las
particiones $\lambda \vdash n$: dos permutaciones son conjugadas si y solo si
tienen la misma estructura de ciclos. Las representaciones irreducibles se
construyen combinatorialmente a partir de \emph{diagramas de Young}.

\begin{definition}[Diagrama de Young y tabla est\'andar]
\label{def:groups:young}
Un \emph{diagrama de Young} de forma $\lambda \vdash n$ es un arreglo de $n$
casillas en filas justificadas a la izquierda, con $\lambda_i$ casillas en la
fila $i$. Una \emph{tabla de Young est\'andar} rellena las casillas con $1,
\dots, n$ de modo que las entradas crecen a lo largo de cada fila y hacia abajo
en cada columna.
\end{definition}

\begin{theorem}[Representaciones irreducibles de $\Sym_n$]
\label{thm:groups:sn-irreps}
Para cada partici\'on $\lambda \vdash n$ hay una representaci\'on irreducible
$V_\lambda$ de $\Sym_n$ cuya dimensi\'on es igual al n\'umero de tablas de Young
est\'andar de forma $\lambda$. Toda representaci\'on irreducible surge de este
modo, y $\sum_{\lambda \vdash n} (\dim V_\lambda)^2 = n!$.
\end{theorem}

\begin{proof}[Construcci\'on]
Para cada tabla $T$ de forma $\lambda$, el politabloide $e_T = a_T b_T$
antisimetriza sobre columnas y simetriza sobre filas; $V_\lambda$ es el subespacio
generado por todos los politabloides. El argumento completo usa el \'algebra de
grupo $\C[\Sym_n]$ y el lema de Schur; v\'eanse \citet{sagan2001symmetric} y
\citet{fulton1997young}.
\end{proof}

Para $\Sym_3$ las particiones y dimensiones son $(3)$ (trivial, $\dim 1$),
$(2,1)$ (est\'andar, $\dim 2$), y $(1,1,1)$ (signo, $\dim 1$), con $1^2 + 2^2 +
1^2 = 6 = 3!$. Aunque elegante, la convoluci\'on expl\'icita sobre $\Sym_n$ v\'ia
representaciones irreducibles es computacionalmente intratable para $n$ grande,
ya que $|\Sym_n| = n!$. Las arquitecturas pr\'acticas evitan trabajar con
$\Sym_n$ directamente y, en su lugar, construyen funciones invariantes y
equivariantes de forma elemental.

\paragraph{DeepSets: aproximaci\'on universal.}
El resultado central para las arquitecturas invariantes a permutaciones
caracteriza por completo la forma de las funciones invariantes continuas.

\begin{theorem}[Aproximaci\'on universal de DeepSets]
\label{thm:groups:deepsets}
Sea $f: \mathcal{X}^n \to \mathcal{Y}$ una funci\'on continua invariante a
permutaciones con $\mathcal{X} \subset \R^d$ e $\mathcal{Y} \subset \R^k$
compactos. Para todo $\epsilon > 0$ existen aplicaciones continuas $\phi:
\mathcal{X} \to \R^m$ y $\rho: \R^m \to \mathcal{Y}$ (con $m$ suficientemente
grande) tales que
\begin{equation}
  \sup_{X \in \mathcal{X}^n}
    \norm{f(X) - \rho\Big(\sum_{i=1}^n \phi(x_i)\Big)} < \epsilon.
  \label{eq:groups:deepsets}
\end{equation}
Adem\'as, $\phi$ y $\rho$ pueden realizarse mediante redes neuronales
\citep{zaheer2017deepsets}.
\end{theorem}

\begin{proof}
El argumento tiene tres pasos. \emph{Paso 1 (reducci\'on a medidas):}
identif\'iquese el conjunto no ordenado $\{x_1, \dots, x_n\}$ con la medida
emp\'irica $\mu_X = \tfrac{1}{n}\sum_i \delta_{x_i}$, de modo que $f$ induce una
funci\'on $\tilde f$ sobre medidas de probabilidad discretas. \emph{Paso 2
(momentos):} por un teorema de tipo Weierstrass sobre espacios de medidas, toda
funci\'on continua sobre medidas de probabilidad se aproxima por $\tilde f(\mu)
\approx F(\int g_1\, d\mu, \dots, \int g_M\, d\mu)$ para funciones base $g_1,
\dots, g_M$ y una $F$ continua; para la medida emp\'irica, $\int g_j\, d\mu_X =
\tfrac{1}{n}\sum_i g_j(x_i)$. \emph{Paso 3 (construcci\'on):} t\'omese $\phi(x) =
(g_1(x), \dots, g_M(x))$ y $\rho(z_1, \dots, z_M) = F(z_1/n, \dots, z_M/n)$.
Entonces $\rho(\sum_i \phi(x_i)) = F(\tfrac1n\sum_i g_1(x_i), \dots) \approx
\tilde f(\mu_X) = f(X)$. El teorema de aproximaci\'on universal para redes
neuronales permite realizar tanto $\phi$ como $\rho$ mediante redes.
\end{proof}

\begin{remark}[Dimensi\'on latente]
\label{rem:groups:latent-dim}
\citet{zaheer2017deepsets} tambi\'en demuestran una representaci\'on exacta $f(X)
= \rho(\sum_i \phi(x_i))$ para un universo numerable, pero con $\phi, \rho$
posiblemente discontinuas. \citet{wagstaff2019limitations} mostraron que, para
una representaci\'on exacta con funciones continuas, la dimensi\'on latente debe
satisfacer $m \geq n$. La suma es el \'unico agregador que garantiza la
aproximaci\'on universal: $\max$ y $\text{mean}$ tienen un poder expresivo
estrictamente menor. Por ejemplo, $\max$ no puede distinguir $\{(1,0),(0,1)\}$ de
$\{(1,1)\}$ bajo $\phi(x) = x$.
\end{remark}

\paragraph{PointNet.}
PointNet \citep{qi2017pointnet} aporta invarianza a permutaciones a las nubes de
puntos tridimensionales, procesando $\mathcal{P} = \{\mathbf{p}_1, \dots,
\mathbf{p}_n\} \subset \R^3$ mediante un codificador compartido por punto, un
pooling sim\'etrico, y un decodificador global:
\begin{equation}
  f_{\text{PointNet}}(\mathcal{P})
  = \rho_\psi\Big(\max_{i=1,\dots,n} \phi_\theta(T \mathbf{p}_i)\Big),
  \label{eq:groups:pointnet}
\end{equation}
donde $T$ es una transformaci\'on de entrada predicha por una T-Net.

\begin{proposition}[Invarianza de PointNet]
\label{prop:groups:pointnet-invariance}
La aplicaci\'on \Cref{eq:groups:pointnet} es exactamente invariante a
permutaciones.
\end{proposition}

\begin{proof}
Para $\pi \in \Sym_n$,
\begin{align}
  f_{\text{PointNet}}(\pi \cdot \mathcal{P})
  = \rho_\psi\Big(\max_i \phi_\theta(T \mathbf{p}_{\pi(i)})\Big)
  = \rho_\psi\Big(\max_j \phi_\theta(T \mathbf{p}_j)\Big)
  = f_{\text{PointNet}}(\mathcal{P}),
\end{align}
puesto que el m\'aximo es una operaci\'on sim\'etrica.
\end{proof}

A diferencia de DeepSets, PointNet usa max-pooling en lugar de suma: aunque el
m\'aximo tiene menor poder expresivo te\'orico
(\Cref{rem:groups:latent-dim}), funciona bien en nubes de puntos donde la
geometr\'ia extremal es informativa. Su limitaci\'on ---tratar cada punto de
forma independiente hasta el pooling global, descartando la estructura local---
motiv\'o PointNet++ \citep{qi2017pointnetplus}, que introduce procesamiento local
jer\'arquico preservando la invarianza a permutaciones.

\paragraph{Set Transformers.}
La agregaci\'on global por suma o m\'aximo descarta informaci\'on sobre las
relaciones entre elementos. Los Set Transformers introducen atenci\'on para que
los elementos interact\'uen preservando la equivarianza a permutaciones. El
bloque de atenci\'on multicabeza sobre un conjunto $X \in \R^{n \times d}$ es
\begin{equation}
  \text{MAB}(X) = \operatorname{Concat}(\text{head}_1, \dots, \text{head}_h) W^O,
  \quad
  \text{head}_i = \Attn(X W_i^Q, X W_i^K, X W_i^V),
\end{equation}
con $\Attn(Q, K, V) = \softmax(QK^\top / \sqrt{d_k}) V$.

\begin{proposition}[Equivarianza de la atenci\'on]
\label{prop:groups:attention-equivariance}
El bloque $\text{MAB}$ es equivariante a permutaciones: $\text{MAB}(\pi \cdot X)
= \pi \cdot \text{MAB}(X)$ para todo $\pi \in \Sym_n$.
\end{proposition}

\begin{proof}
El producto $QK^\top$ y la posterior multiplicaci\'on por $V$ permutan las filas
de forma coherente con la permutaci\'on de $X$, mientras que softmax act\'ua de
forma independiente sobre cada fila, preservando la permutaci\'on. La
composici\'on es, por tanto, equivariante.
\end{proof}

Para controlar el coste, el Induced Set Attention Block usa $m$ puntos inductores
aprendibles $I \in \R^{m \times d}$, reduciendo la complejidad de $O(n^2)$ a
$O(nm)$ con $m \ll n$:
\begin{equation}
  \text{ISAB}_m(X) = \text{MAB}(X, \text{MAB}(I, X)),
\end{equation}
donde $\text{MAB}(A, B)$ denota la atenci\'on con consultas de $A$ y
claves/valores de $B$. El Pooling by Multihead Attention usa $k$ semillas
aprendibles $S \in \R^{k \times d}$ para producir una salida de tama\~no fijo e
invariante a permutaciones: $\text{PMA}_k(X) = \text{MAB}(S, \text{rFF}(X))$, con
$\text{rFF}$ una capa feedforward por filas. Estos mecanismos de atenci\'on sobre
conjuntos son la versi\'on equivariante a permutaciones de la atenci\'on
introducida en \Cref{ch:grids}.

\begin{example}[DeepSets sobre un conjunto peque\~no]
\label{ex:groups:deepsets-numerical}
T\'omese $\mathcal{X} = \{(2,1),(1,3),(3,2)\}$ con el codificador lineal
$\phi(\mathbf{x}) = W\mathbf{x} + \mathbf{b}$, $W = \left(\begin{smallmatrix} 1 &
-1 \\ 0 & 2 \\ 1 & 1 \end{smallmatrix}\right)$, $\mathbf{b} = (0,1,0)^\top$.
Entonces $\phi((2,1)) = (1,3,3)^\top$, $\phi((1,3)) = (-2,7,4)^\top$,
$\phi((3,2)) = (1,5,5)^\top$, y la agregaci\'on sim\'etrica da $\sum_i
\phi(\mathbf{x}_i) = (0,15,12)^\top$. Con el decodificador $\rho(\mathbf{z}) =
V\mathbf{z}$, $V = \left(\begin{smallmatrix} 1 & 0 & 0 \\ 0 & 0.1 & 0
\end{smallmatrix}\right)$, la salida es $f(\mathcal{X}) = (0, 1.5)^\top$. Permutar
el conjunto deja la suma ---y por tanto $f$--- sin cambios, confirmando la
invarianza a permutaciones.
\end{example}

Las arquitecturas invariantes y equivariantes a permutaciones han transformado
dominios donde los datos carecen de un orden natural: clasificaci\'on,
segmentaci\'on y detecci\'on sobre nubes de puntos LiDAR; predicci\'on y
generaci\'on de propiedades moleculares, donde el etiquetado de los \'atomos es
arbitrario; identificaci\'on de chorros y detecci\'on de anomal\'ias en f\'isica
de altas energ\'ias, donde el orden de las part\'iculas detectadas no porta
informaci\'on; y razonamiento abstracto sobre conjuntos. Son tambi\'en el
fundamento conceptual de las redes neuronales de grafos del \Cref{ch:graphs},
donde la invarianza al reetiquetado de nodos es la simetr\'ia definitoria.

\section{Atenci\'on equivariante}
\label{sec:groups:equivariant-attention}

El mismo principio que produce la convoluci\'on de grupo produce tambi\'en
\emph{atenci\'on equivariante}: incorporar simetr\'ias geom\'etricas al mecanismo
de atenci\'on de modo que los pesos adaptativos, dependientes del contenido,
respeten la estructura de los datos. Adaptamos la atenci\'on del \Cref{ch:grids}
para que act\'ue equivariantemente bajo un grupo $\group$.

\begin{theorem}[Construcci\'on de la atenci\'on equivariante]
\label{thm:groups:equivariant-attention}
Para un grupo $\group$ que act\'ua sobre el espacio de caracter\'isticas, un
mecanismo de atenci\'on es $\group$-equivariante si y solo si (i) las funciones
de compatibilidad son $\group$-invariantes y (ii) las transformaciones de
consulta, clave y valor conmutan con la acci\'on del grupo.
\end{theorem}

\begin{proof}[Esbozo]
($\Leftarrow$) Si la compatibilidad $a(\mathbf{q}_i, \mathbf{k}_j)$ es
invariante, entonces $a(\rho(g)\mathbf{q}_i, \rho(g)\mathbf{k}_j) =
a(\mathbf{q}_i, \mathbf{k}_j)$, de modo que los pesos de atenci\'on $\alpha_{ij}$
son invariantes; si adem\'as $W_V$ conmuta con $\group$ ($W_V \rho(g) = \rho(g)
W_V$), entonces $\sum_j \alpha_{ij} W_V \rho(g)\mathbf{x}_j = \rho(g) \sum_j
\alpha_{ij} W_V \mathbf{x}_j$, lo cual es equivarianza. ($\Rightarrow$) Si la
atenci\'on es equivariante, los pesos deben ser invariantes bajo la acci\'on
simult\'anea sobre consultas y claves, forzando que las funciones de
compatibilidad sean invariantes; la equivarianza de la salida con pesos fijos
fuerza entonces que $W_V$ conmute con $\rho(g)$.
\end{proof}

\begin{lemma}[Condici\'on suficiente para la equivarianza de la autoatenci\'on]
\label{lem:groups:self-attention-equivariance}
Sea $\group$ que act\'ua sobre $\R^{n \times d}$ por permutaciones. Si las
matrices de proyecci\'on satisfacen $\mathbf{W}_Q = \mathbf{W}_K = \mathbf{W}_V =
\mathbf{W}$ y $\mathbf{W}$ conmuta con la acci\'on del grupo, entonces la
autoatenci\'on es $\group$-equivariante.
\end{lemma}

\begin{proof}
Para una matriz de permutaci\'on $P \in \group$ y una entrada $\mathbf{X}$,
\begin{align}
  \text{SelfAttention}(P\mathbf{X})
  &= \softmax\!\Big(\tfrac{(P\mathbf{X}\mathbf{W})(P\mathbf{X}\mathbf{W})^\top}
     {\sqrt{d_k}}\Big)(P\mathbf{X}\mathbf{W}) \\
  &= \softmax\!\Big(\tfrac{P(\mathbf{X}\mathbf{W})(\mathbf{X}\mathbf{W})^\top
     P^\top}{\sqrt{d_k}}\Big)P(\mathbf{X}\mathbf{W})
   = P \cdot \text{SelfAttention}(\mathbf{X}),
\end{align}
donde el \'ultimo paso usa la conmutatividad y la equivarianza de softmax a las
permutaciones.
\end{proof}

\paragraph{Atenci\'on irreducible v\'ia Clebsch--Gordan.}
Para grupos de Lie compactos, la teor\'ia de representaciones de
\Cref{sec:groups:tfn} subyace a la atenci\'on equivariante. Las caracter\'isticas
de cada capa se descomponen por tipo de representaci\'on, $\mathbf{f} =
\bigoplus_{\ell=0}^{L_{\max}} \mathbf{f}^{(\ell)}$ con $\mathbf{f}^{(\ell)} \in
\R^{n_\ell \times (2\ell+1)}$, y los arm\'onicos esf\'ericos $Y_\ell^m$ realizan
las representaciones irreducibles $(2\ell+1)$-dimensionales de $\SO(3)$,
transformando bajo las matrices de Wigner $D^{(\ell)}(R)$
\citep{brocker1985representations}. Combinar caracter\'isticas de distintos tipos
requiere el producto de Clebsch--Gordan (\Cref{thm:groups:clebsch-gordan}); por
ejemplo $D^{(1)} \otimes D^{(1)} = D^{(0)} \oplus D^{(1)} \oplus D^{(2)}$ produce
un escalar (producto escalar), un pseudovector (producto vectorial), y un tensor
sim\'etrico sin traza. Una puntuaci\'on de atenci\'on equivariante combina
consultas y claves de distintos tipos a trav\'es de estos productos,
\begin{equation}
  e_{ij} = \sum_{\ell_q, \ell_k, \ell_{\text{out}}}
    \MLP_{\ell_q, \ell_k, \ell_{\text{out}}}
    \big(\mathbf{q}_i^{(\ell_q)} \otimes \mathbf{k}_j^{(\ell_k)}
    \big)^{(\ell_{\text{out}})},
\end{equation}
con los perceptrones multicapa actuando de forma independiente dentro de cada
tipo de representaci\'on.

\paragraph{Arquitecturas de transformer equivariante.}
Varias arquitecturas instancian este marco. Los $\SE(3)$-Transformers de
\citet{fuchs2020se3transformers}, ya analizados en \Cref{sec:groups:tfn},
calculan pesos de atenci\'on invariantes a partir de caracter\'isticas escalares,
distancias, y los productos internos $\mathbf{f}_{1,i} \cdot \mathbf{f}_{1,j}$ de
las caracter\'isticas vectoriales, y luego combinan mensajes de valor
equivariantes; la equivarianza se sigue exactamente como en
\Cref{thm:groups:se3-transformer-equivariance}. El Equiformer de
\citet{liao2022equiformer} reemplaza la puntuaci\'on de producto escalar
$\mathbf{q}_i^\top \mathbf{k}_j$ por una atenci\'on de perceptr\'on multicapa
equivariante $e_{ij} = \MLP_{\text{attn}}(\mathbf{q}_i \otimes \mathbf{k}_j,
\norm{\mathbf{r}_{ij}})$, donde $\otimes$ es el producto tensorial de
caracter\'isticas irreducibles, permitiendo interacciones m\'as ricas que una
\'unica puntuaci\'on escalar a la vez que preserva la equivarianza $\SO(3)$. El
LieTransformer de \citet{hutchinson2020lietransformer} proporciona un marco
general para grupos de Lie arbitrarios elevando las caracter\'isticas al grupo,
aplicando atenci\'on en el espacio del grupo, y proyectando de vuelta,
\begin{equation}
  \text{LieAttn}(\mathbf{x}) = \text{Project}\big(
    \Attn(\text{Lift}(\mathbf{x}))\big),
\end{equation}
lo que garantiza la equivarianza a trav\'es de la representaci\'on regular
izquierda $(L_g f)(h) = f(g^{-1} h)$. Extendiendo el principio al
preentrenamiento multidominio, el Equivariant Pretrained Transformer de
\citet{jiao2024ept} es un modelo fundacional que unifica mol\'eculas peque\~nas y
prote\'inas bajo un \'unico marco $\Eucl(3)$-equivariante, explotando que la
equivarianza es una propiedad arquitect\'onica, preservada bajo la transferencia
a cualquier dominio que admita la misma acci\'on de grupo.

\paragraph{Atenci\'on vectorial sobre nubes de puntos.}
Una l\'inea complementaria intercambia la equivarianza $\SO(3)$ exacta por
eficiencia. El Point Transformer de \citet{zhao2021point} usa atenci\'on
\emph{vectorial}, en la que el peso de atenci\'on es un vector que modula los
canales de caracter\'isticas de forma independiente:
\begin{equation}
  \mathbf{y}_i = \sum_{j \in \mathcal{N}(i)}
    \rho\big(\gamma(\phi(\mathbf{x}_i) - \psi(\mathbf{x}_j) + \delta)\big)
    \big(\alpha(\mathbf{x}_j) + \delta\big),
\end{equation}
donde $\phi, \psi, \alpha$ son aplicaciones lineales, $\gamma, \rho$ son
perceptrones multicapa, y $\delta = \MLP(\mathbf{p}_i - \mathbf{p}_j)$ codifica
la posici\'on relativa. Como $\delta$ depende solo de la posici\'on relativa
$\mathbf{p}_i - \mathbf{p}_j$, los pesos de atenci\'on son invariantes a
traslaciones, conectando directamente con la condici\'on de invarianza de
\Cref{thm:groups:equivariant-attention}. Evitando los productos de
Clebsch--Gordan de coste $O(\ell_{\max}^3)$, el Point Transformer alcanza una
complejidad $O(|\mathcal{N}(i)| \cdot d)$ por punto, al precio de renunciar a la
equivarianza $\SO(3)$ exacta de las caracter\'isticas de salida.

\section{Una visi\'on unificada}
\label{sec:groups:unified}

Las arquitecturas de este cap\'itulo realizan un \'unico principio en escenarios
cada vez m\'as generales. La elecci\'on del grupo de simetr\'ia determina tanto la
capacidad representacional como el coste computacional, resumidos en
\Cref{tab:groups:comparison}.

\begin{table}[h]
\centering
\small
\caption{Grupos de simetr\'ia, complejidad y comportamiento de las arquitecturas
equivariantes. $N$: resoluci\'on espacial; $L$: grado esf\'erico m\'aximo; $F$:
frecuencia arm\'onica m\'axima.}
\label{tab:groups:comparison}
\begin{tabular}{@{}lcccl@{}}
\toprule
\textbf{Grupo} & \textbf{Orden} & \textbf{Tipo} & \textbf{Complejidad} & \textbf{Dominio} \\
\midrule
CNN est\'andar & --- & Traslacional & $O(N^2)$ & Rejillas \\
$C_4$ & $4$ & Finito & $O(4 N^2)$ & Rotaciones discretas \\
$D_4$ & $8$ & Finito & $O(8 N^2)$ & Rotaciones + reflexiones \\
$\SO(2)$ orientable & $\infty$ & Continuo & $O(F^2 N^2)$ & Rotaciones planas continuas \\
$\SO(3)$ esf\'erico & $\infty$ & Continuo & $O(L^3 N^2)$ & Datos esf\'ericos \\
$\SE(3)$ & $\infty$ & No compacto & $O(L^4 N^2)$ & Nubes de puntos, mol\'eculas \\
$\Sym_n$ & $n!$ & Finito & $O(n^2)$ & Conjuntos, grafos \\
\bottomrule
\end{tabular}
\end{table}

Del compromiso entre fidelidad de la equivarianza y coste computacional se siguen
tres directrices. Para escenarios en tiempo real, las redes convolucionales de
grupo discreto ($C_4$, $D_4$) dan el mejor equilibrio entre ganancia de precisi\'on
y sobrecarga. Para alta precisi\'on, las arquitecturas orientables y $\SE(3)$
justifican su mayor coste con una precisi\'on superior. Cuando las simetr\'ias se
desconocen, conviene empezar con redes convolucionales de grupo simples para
confirmar su presencia antes de adoptar maquinaria m\'as pesada, siguiendo el
camino iterativo: red convolucional est\'andar $\to$ red convolucional de grupo
discreto $\to$ extensiones continuas.

El hilo \'unico que recorre el cap\'itulo es que la equivarianza convierte una
simetr\'ia de los datos en una restricci\'on sobre la arquitectura, y esa
restricci\'on produce eficiencia en par\'ametros, generalizaci\'on robusta, y una
receta constructiva para redes completas. Las convoluciones de grupo
desarrolladas aqu\'i generalizan la convoluci\'on cl\'asica de las traslaciones a
grupos de simetr\'ia arbitrarios, pero a\'un suponen simetr\'ias algebraicas
\emph{globales} y exactas. El siguiente paso en la jerarqu\'ia del aprendizaje
profundo geom\'etrico es manejar datos sobre espacios geom\'etricos generales,
donde las simetr\'ias son locales o solo aproximadas. Descartar por completo la
rejilla ambiente y conservar \'unicamente la estructura relacional conduce a las
redes neuronales de grafos del \Cref{ch:graphs}; dotar al dominio de una distancia
y una curvatura intr\'insecas conduce a las variedades y geod\'esicas del
\Cref{ch:geodesics}; y reconocer que un dominio curvo no admite un marco local
can\'onico conduce a las redes equivariantes gauge del \Cref{ch:gauges}.
